\documentclass[11pt,oneside]{book}
\usepackage[margin=1.2in]{geometry}
\usepackage[toc,page]{appendix}
\usepackage{graphicx}
\usepackage{natbib}
\usepackage{lipsum}
\usepackage{caption}
\usepackage[T1]{fontenc}
\usepackage{titlesec, blindtext, color}
\usepackage{xcolor,tikz}
\usetikzlibrary{patterns}
\usetikzlibrary{decorations.markings}
\usepackage{amsmath,amssymb,amsthm,mathrsfs,amsfonts,xfrac,pifont,bbold,physics}
\usepackage[utf8]{inputenc}
\usepackage{amsthm}
\usepackage[breakable, theorems, skins]{tcolorbox}
\usepackage[colorlinks = true,
            linkcolor = red,
            urlcolor  = blue,
            citecolor = red,
            anchorcolor = red]{hyperref}
\usepackage{cleveref}

\usepackage{soul}
\usepackage{frcursive}
\usepackage{booktabs}
\usepackage{enumitem}

\usepackage{array}
\newcolumntype{C}[1]{>{\centering\arraybackslash}m{#1}}

% -------------------------------------------------------------------
% Theorem Styles
% -------------------------------------------------------------------

\theoremstyle{definition} % Define theorem styles here based on the definition style (used for definitions and examples)
\newtheorem*{definition}{Definition}

\theoremstyle{plain} % Define theorem styles here based on the plain style (used for theorems, lemmas, propositions)
\newtheorem{theorem}{Theorem}[section]
\newtheorem{axiom}{Axiom}
\newtheorem{corollary}[theorem]{Corollary}
\newtheorem{lemma}[theorem]{Lemma}
\newtheorem{proposition}[theorem]{Proposition}

\theoremstyle{remark} % Define theorem styles here based on the remark style (used for remarks and notes)
\newtheorem*{notation}{Notation}
\newtheorem*{solution}{Solution}

\newtheoremstyle{underline}% name
{}        % Space above, empty = `usual value'
{}              % Space below
{}              % Body font
{}    % Indent amount (empty = no indent, \parindent = para indent)
{}              % Thm head font
{.}             % Punctuation after thm head
{1.5mm}         % Space after thm head: \newline = linebreak
{{\underline{\textit{\thmname{#1}\thmnumber{ #2}}~\thmnote{(#3)}\unskip}}}  % Thm head spec

\theoremstyle{underline}

\newtheorem{remark}[theorem]{Remark}
\newtheorem{example}[theorem]{Example}
\newtheorem{claim}[theorem]{Claim}
% -------------------------------------------------------------------
% Chapter Headings
% -------------------------------------------------------------------

\setcounter{chapter}{-1}

\makeatletter
\renewcommand{\@chapapp}{Chapter}
\makeatother
\definecolor{lightergray}{rgb}{0.9,0.9,0.9}

\usepackage{titlesec}
\titleformat{\section}{\large\bfseries\raggedright}{}{0em}{\colorsection}[\titlerule]
\titleformat{name=\section,numberless}{\large\scshape\bfseries\raggedright}{}{0em}{\colorsectionnonumber}[\titlerule]

\titleformat{\subsection}{\bfseries\raggedright}{}{0em}{\colorsubsection}
\titleformat{name=\subsection,numberless}{\bfseries\raggedright}{}{0em}{\colorsubsectionnonumber}

\newcommand{\colorsection}[1]{%
    \colorbox{lightergray}{\parbox{\dimexpr\textwidth-2\fboxsep}{\thesection\ \ #1}}}
\newcommand{\colorsectionnonumber}[1]{%
    \colorbox{lightergray}{\parbox{\dimexpr\textwidth-2\fboxsep}{#1}}}
    
\newcommand{\colorsubsection}[1]{%
    \colorbox{lightergray}{\parbox{\dimexpr\textwidth-2\fboxsep}{\thesubsection\ #1}}}
\newcommand{\colorsubsectionnonumber}[1]{%
    \colorbox{lightergray}{\parbox{\dimexpr\textwidth-2\fboxsep}{#1}}}
    
\definecolor{gray75}{gray}{0.75}
\newcommand{\hsp}{\hspace{20pt}}
\titleformat{\chapter}[hang]{\Huge\bfseries}{\thechapter\hsp\textcolor{gray75}{|}\hsp}{0pt}{\Huge\bfseries}

% Begin/end Shortcuts
\def\bc{\begin{corollary}}
\def\ec{\end{corollary}}
\def\bcl{\begin{claim}}
\def\ecl{\end{claim}}
\def\bd{\begin{definition}}
\def\ed{\end{definition}}
\def\ben{\begin{enumerate}}
\def\een{\end{enumerate}}
\def\be{\begin{equation}}
\def\ee{\end{equation}}
\def\bse{\begin{equation*}}
\def\ese{\end{equation*}}
\def\bex{\begin{example}}
\def\eex{\end{example}}
\def\bit{\begin{itemize}}
\def\eit{\end{itemize}}
\def\bl{\begin{lemma}}
\def\el{\end{lemma}}
\def\bnn{\begin{notation}}
\def\enn{\end{notation}}
\def\bn{\begin{note}}
\def\en{\end{note}}
\def\bp{\begin{proposition}}
\def\ep{\end{proposition}}
\def\bq{\begin{proof}}
\def\eq{\end{proof}}
\def\br{\begin{remark}}
\def\er{\end{remark}}
\def\bs{\begin{solution}}
\def\es{\end{solution}}
\def\btab{\begin{table}}
\def\etab{\end{table}}
\def\btb{\begin{tabular}}
\def\etb{\end{tabular}}
\def\bt{\begin{theorem}}
\def\et{\end{theorem}}
\def\btik{\begin{tikzpicture}}
\def\etik{\end{tikzpicture}}

% Letter Shortcuts
% Greek letters
\def\a{\alpha}
\def\b{\beta}
\def\del{\delta}
\def\g{\gamma}
\def\l{\lambda}
\def\n{\nabla}
\def\sig{\sigma}

% Mathbb
\def\b1{\mathbb{1}} % identity
\def\C{\mathbb{C}} % complex numbers
\def\F{\mathbb{F}}
\def\N{\mathbb{N}} % Natural
\def\Q{\mathbb{Q}} % Rational
\def\R{\mathbb{R}} % Reals
\def\Z{\mathbb{Z}} % Integers

% Mathcal
\def\cA{\mathcal{A}}
\def\cB{\mathcal{B}}
\def\cD{\mathcal{D}}
\def\cF{\mathcal{F}}
\def\cH{\mathcal{H}}
\def\cJ{\mathcal{J}}
\def\cK{\mathcal{K}}
\def\cL{\mathcal{L}}
\def\cM{\mathcal{M}}
\def\cN{\mathcal{N}}
\def\cO{\mathcal{O}}
\def\cP{\mathcal{P}}
\def\cT{\mathcal{T}}

% Mathfrak 
\DeclareMathOperator{\GL}{GL}
\def\gl{\mathfrak{gl}}
\DeclareMathOperator{\Ort}{O}
\def\ort{\mathfrak{o}}
\DeclareMathOperator{\SL}{SL}
\def\sl{\mathfrak{sl}}
\DeclareMathOperator{\SO}{SO}
\def\so{\mathfrak{so}}
\DeclareMathOperator{\SU}{SU}
\def\su{\mathfrak{su}}
\def\fF{\mathfrak{F}}

% Misc. 
\def\diag{\text{diag}}
\def\p{\partial}
\def\ss{\subset}
\def\se{\subseteq}
\def\sm{\setminus}
\def\ve{\varepsilon}
\def\vn{\varnothing}
\def\cl{\colon}
\def\la{\langle}
\def\ra{\rangle}

% Boxes 
\DeclareRobustCommand{\mybox}[2][colframe=blue!10!black,before skip=10pt]{%
\begin{tcolorbox}[
        breakable,
        width=\dimexpr\textwidth\relax, 
        boxsep=5pt,
        arc=5pt,outer arc=5pt,
        ]
        #2
\end{tcolorbox}
}
\def\bbox{\begin{tcolorbox}[colback=white!5,colframe=blue!75!black,title=Exercise]}
\def\ebox{\end{tcolorbox}}

% Feynman Diagram Stuff 
\def\midarrow{\draw[thick, decoration={markings, mark=at position 0.5 with {\arrow{>}}}, postaction={decorate}]}
\def\midarrowblue{\draw[thick, blue, decoration={markings, mark=at position 0.5 with {\arrow{>}}}, postaction={decorate}]}
\def\midarrowred{\draw[thick, red, decoration={markings, mark=at position 0.5 with {\arrow{>}}}, postaction={decorate}]}
\def\aftermidarrow{\draw[thick, decoration={markings, mark=at position 0.7 with {\arrow{>}}}, postaction={decorate}]}
\def\aftermidarrowblue{\draw[thick, blue, decoration={markings, mark=at position 0.7 with {\arrow{>}}}, postaction={decorate}]}
\def\aftermidarrowred{\draw[thick, red, decoration={markings, mark=at position 0.7 with {\arrow{>}}}, postaction={decorate}]}
\def\beforemidarrow{\draw[thick, decoration={markings, mark=at position 0.4 with {\arrow{>}}}, postaction={decorate}]}
\def\wavey{\draw[thick, decorate, decoration={snake, segment length=2.5mm, amplitude=0.75mm}]}
\newcommand{\AxisRotatorL}[1][rotate=0]{%
    \tikz [x=0.25cm,y=0.60cm,line width=.2ex,-stealth,#1] \draw (0,0) arc (-150:150:1 and 1);%
}
\newcommand{\AxisRotatorR}[1][rotate=0]{%
    \tikz [x=0.25cm,y=0.60cm,line width=.2ex,-stealth,#1] \draw (0,0) arc (150:-150:1 and 1);%
}

\begin{document}

\captionsetup[figure]{margin=1.5cm,font=small,labelfont={bf},name={Figure},labelsep=colon,textfont={it}}
\captionsetup[table]{margin=1.5cm,font=small,labelfont={bf},name={Table},labelsep=colon,textfont={it}}
\setlipsumdefault{1}

\frontmatter

\begin{titlepage}
	\centering
	    \scshape % Use small caps for all text on the title page
        \vspace*{\baselineskip} % White space at the top of the page
        
	    \rule{\textwidth}{1.6pt}\vspace*{-\baselineskip}\vspace*{2pt} % Thick horizontal rule
	    \rule{\textwidth}{0.4pt} % Thin horizontal rule
	    
	    \vspace{0.75\baselineskip} % Whitespace above the title
	    
	    {\LARGE Amplitudes} % Title
	    
	    \vspace{0.75\baselineskip} % Whitespace below the title
	    
	    \rule{\textwidth}{0.4pt}\vspace*{-\baselineskip}\vspace{3.2pt} % Thin horizontal rule
	    \rule{\textwidth}{1.6pt} % Thick horizontal rule
	    
        \vspace{5\baselineskip} % Whitespace after the title block

	    Course delivered in 2020 by 
	
	    \vspace{0.5\baselineskip} % Whitespace before the editors
	
	    {\scshape\Large Dr. Arthur Lipstein} % Lecturer Name
	
	    \vspace{0.5\baselineskip} % Whitespace below the editor list
	
	    \textit{Durham University} % Lecturer Institution
	    
	    \vspace{2\baselineskip} % Whitespace after the title block
	    
	    \includegraphics[width=8cm]{images/DurhamLogo.png}\\[1cm] % Logo 

	    Notes taken by 
	
	    \vspace{0.5\baselineskip} % Whitespace before my name
	
	    {\scshape\Large Richie Dadhley} % Lecturer Name
	   
	    \vspace{0.5\baselineskip} % Whitespace below my name
	    \textit{richie@dadhley.com} % Email
	
	    \vfill % Whitespace between editor names and publisher logo
\end{titlepage}



% -------------------------------------------------------------------
% Acknowledgements
% -------------------------------------------------------------------

\newpage
\section*{Acknowledgements}

These are my notes on the 2020 lecture course "Amplitudes" taught by Dr. Arthur Lipstein at Durham University as part of the Particles, Strings and Cosmology Msc. For reference, the course lasted 8 hours and was taught over 4 weeks. \\

I have tried to correct any typos and/or mistakes I think I have noticed over the course. I have also tried to include additional information that I think supports the taught material well, which sometimes has resulted in modifying the order the material was taught. Obviously, any mistakes made because of either of these points are entirely mine and should not reflect on the taught material in any way. \\

% For any other notable sources used:
% One of the main sources I have used to obtain the additional information is Dr. David Tong’s String Theory notes, available via the following link  

%\begin{center}
%    \textcolor{red}{Link}
%\end{center}

I would like to extend a message of thanks to Dr. Lipstein for teaching this course. \\

If you have any comments and/or questions please feel free to contact me via the email provided on the title page. \\

These notes are now finished. If you have any comments and/or questions please feel free to contact me via the email provided on the title page. \\

For a list of other notes/works I have available, visit my blog site:

\begin{center}
    \href{https://richie291.wixsite.com/theoreticalphysics}{https://richie291.wixsite.com/theoreticalphysics}
\end{center}

These notes are not endorsed by Dr. Lipstein or Durham University.

\vspace{1cm}

\begin{flushright}
    \Huge{{\cursive\setul{0.1ex}{}\ul{~~Richie Dadhley~~}}}
\end{flushright}

% -------------------------------------------------------------------
% Contents
% -------------------------------------------------------------------

\tableofcontents

% -------------------------------------------------------------------
% Main sections (as required)
% -------------------------------------------------------------------

\mainmatter

\chapter{Introduction}

Scattering amplitudes are the bread and butter of modern particle physics. Broadly speaking, they describe the probability for a collection of particles to interact in a certain way. One of the things that has been deriving the research in this area is that they have remarkably simple mathematical constructions. 

We are (hopefully) familiar with amplitudes\footnote{From now on we may simply say "amplitudes" to mean scattering amplitudes.} from our previous QFT courses and we know that we can calculate them from the Feynman diagrams. The natural question that arises, then, is "why do we want a whole other course dedicated to this?" Well, as we will hopefully convince, there is an alternate way to compute amplitudes that drastically reduces the complexity of the computations and makes manifest these simplifications. In order to do that, of course it is first useful to remind ourselves what exactly an amplitude is. 

We will also clear up our conventions in the rest of this introduction. 

\section{What Is A Scattering Amplitude?}

\bd 
    We define a scattering amplitude as a matrix element 
    \be 
    \label{eqn:Amplitude}
        \cA = \bra{\text{out}} S \ket{\text{in}}
    \ee 
    with our states being asymptotic past/future states --- i.e. $\ket{\text{in}}$ is evaluated at $t\to -\infty$ and $\bra{\text{out}}$ at $t\to +\infty$ --- they are given by direct products of single particle states with definite \textit{on-shell} momentum and polarisation. We denote them as follows $\ket{p_i,\epsilon_i}$. 
\ed 

\br  
    We shall work with the "mostly minus", $(+,-,-,-)$, sign convention. 
\er 

As we are aware, it is convenient to split $\cA$ into two parts which correspond to connected and unconnected diagrams. We do this by setting 
\bse
    S = \b1 + i T
\ese 
with $T$ being the \textit{transition amplitude}. Next recall that the LSZ formula states that\footnote{More details about the LSZ reduction formula can be found in, for example, Chapter 8 of my IFT notes or section 3.2 of my QFT II notes.} 
\bse 
    \bra{\text{out}} iT \ket{\text{in}} = \prod_{i=1}^{n_i+n_f} \sqrt{Z_i} \, ( ... )
\ese 
with the $...$ stands for "sum of all amputated and connected Feynman diagrams dressed with polarisations, with $n_i$ incoming states and $n_f$ outgoing states." The factor $Z_i$ is known as the \textit{wavefunction renormalisation} and it can be computed from the residue of the 2-point function. The observable thing in the end is the so-called \textit{differential cross section}, which is obtained from $|T|^2$. 

\section{What's Wrong With The Feynman Diagram Approach?}

We now return to the comment we made at the start of this introduction that although in principle we know how to compute amplitudes from the Feynman diagrams that this course aims to present an alternate method quoting the mantra of "it simplifies the calculations". Let's expand a little bit on why this is the case now.

Of course Feynman diagrams are an extremely useful tool to particle physicists, but nothing is perfect and they do indeed have problems. This is not some new profound statement but simply corresponds to the fact that when we want to work to higher accuracy with Feynman diagrams we necessarily have to include more diagrams. As well as this calculating more complicated processes naturally involves considering more particles and therefore more terms. It is easy to see that we end up needing to process \textit{huge} amounts of information. For example, the $n$-point gluon amplitude grows as $n!$.

The big motivational point to notice is that, once these heavy computations are done, we end up obtaining extremely nice expressions. That is we start of with some highly complicated collection of data but cranking the handle on the calculation leads to huge simplifications at some point. The idea of this course is to try reformulate the calculations in order to make this simplifications manifest. 

To give one last guiding remark, one of the main problems with the Feynman diagram approach we seek to counteract is the fact that individual diagrams are \textit{not} gauge invariant; it is only their \textit{sum} that is. Feynman diagrams also involve the computation of \textit{off-shell}, virtual, particles. However, as we should be aware, the amplitude must be gauge invariant (it's physically observable!) and it knows nothing about off-shell degrees of freedom (our states in \Cref{eqn:Amplitude} are \textit{on-shell}). 

The ultimate message is that the methods presented in this course are more than just being more efficient/being clever, but they point towards new, alternative reformulations of physics. 

\br 
    A warning before proceeding: A lot of the initial work is going to be setting up efficient notation. This might be boring at first, but it will have huge pay offs later. So bare with it. 
\er 

\br 
    Ultimately\footnote{At least for $\cN=4$ Super Yang-Mills theory.} the reformulation presented here involves elements of twistor string theory. Of course we will not go into the details here, but this remark is just to highlight this fact. 
\er
\chapter{Spinor-Helicity Formalism}

As we just said in the introduction, the idea is to reformulate stuff using the fact that we know that the amplitudes are gauge invariant and only know about on-shell degrees of freedom. The spinor-helicity formalism does just this: it is meant to make the on-shell-ness of the amplitudes manifest, i.e. the momentum is on shell and the polarisation are gauge invariant. 

Consider a 4d \textit{on-shell} momentum. We can rewrite it in a slightly funny way
\be 
\label{eqn:pAlphaDotAlpha}
    p^{\dot{\a}\a} := p_{\mu} (\bar{\sig}^{\mu})^{\dot{\a}\a}  = \begin{pmatrix}
        p^0+p^3 & p^1-ip^2 \\
        p^1+ip^2 & p^0-p^3
    \end{pmatrix}, 
\ee 
where 
\bse 
    \bar{\sig}^{\mu} = (\b1, -\vec{\sig}) 
\ese 
and $\dot{\a},\a=1,2$. This should be familiar from the SUSY course, but to refresh our memories is that locally\footnote{I.e. the Lie algebra.} the Lorentz group is the same as $SU(2)\times SU(2)$. The indices $\a, \dot{\a}$ then correspond to the fundamental representations in the two $SU(2)$s. Elements in each of the $SU(2)$s are known as \textit{spinors}. 

The obvious question is "why do we write the momentum in this seemingly more complicated notation?" The answer comes from considering the determinant of \Cref{eqn:pAlphaDotAlpha}:
\bse 
    \det (p^{\dot{\a}\a}) = (p^0)^2 - (\vec{p})^2 = p^2 = m^2,
\ese 
where the last line follows from our assumption that the momentum is on-shell. So the determinant gives us the mass (squared), what use is this? Well now consider the case when $m=0$. Obviously this tells us that the matrix is degenerate, which in turn tells us that $p^{\dot{\a}\a}$ has rank-$1$. This means that we can decompose $p^{\dot{\a}\a}$ as the outer product\footnote{The outer product of two matrices is simply their tensor product. For our case of 2-column vectors, $\l^{\a} = (\l^1,\l^2)^T$ and $\widetilde{\l}^{\dot{\a}} = (\widetilde{\l}^{\dot{1}},\widetilde{\l}^{\dot{2}})^T$, it is simply given by \bse 
    \widetilde{\l}^{\dot{\a}}\otimes \l^{\a} = \begin{pmatrix}
        \widetilde{\l}^{\dot{1}} \l^1 & \widetilde{\l}^{\dot{1}}\l^2 \\
        \widetilde{\l}^{\dot{2}} \l^1 & \widetilde{\l}^{\dot{2}}\l^2 \\
    \end{pmatrix}.
\ese 
We simply suppress the $\otimes$ in \Cref{eqn:BispinorForm} and will continue to do so throughout these notes. Hopefully it will be clear from context ($p^{\dot{\a}\a}$ is a $2\times 2$ matrix whereas the $\widetilde{\l}/\l$ are 2-column matrices) what is meant.} of two spinors:
\mybox{
    \be
    \label{eqn:BispinorForm}
        p^{\dot{\a}\a} = \widetilde{\l}^{\dot{\a}} \l^{\a}
    \ee 
}
\noindent which is known as \textit{bispinor form}. 

\br
\label{rem:CommutingSpinors}
    In these notes we will work with \textit{commuting} spinors. That is, for a given $\a/\dot{\a}$, our $\l^{\a}$ and $\widetilde{\l}^{\dot{\a}}$ are `normal', Grassman even, numbers that we can freely interchange. This is in contrast to our convention in SUSY where we had \textit{anticommuting} spinors. The reason we can do this is not trivial at this point, and will be explained later. The idea is that we split the amplitude into a commuting part and an anticommuting part. The latter part is the colour structure and is contained completely within the generators $T^a$. 
\er 

\br 
    Note that in the above remark we said "for a given $\a/\dot{\a}$" because it is only the \textit{components} we can commute. In other words
    \bse 
        \l^{\a} \widetilde{\l}^{\dot{\a}} := \l^{\a} \otimes \widetilde{\l}^{\dot{\a}} \cong \widetilde{\l}^{\dot{\a}} \otimes \l^{\a} =: \widetilde{\l}^{\dot{\a}}\l^{\a},
    \ese 
    where $\cong$ means "isomorphic as algebras", but they are not \textit{equal}. This is easily seen by the act that $p^{\a\dot{\a}} \neq p^{\dot{\a}\a}$ using \Cref{eqn:pAlphaDotAlpha}.
\er 

\bbox 
    Let 
    \bse 
        \l^{\a} = \frac{1}{\sqrt{p^0+p^3}} \begin{pmatrix}
            p^0 + p^3 \\
            p^1 - ip^2
        \end{pmatrix}, \qand \widetilde{\l}^{\dot{\a}} = \frac{1}{\sqrt{p^0+p^3}} \begin{pmatrix}
            p^0 + p^3 \\
            p^1 + ip^2
        \end{pmatrix}.
    \ese 
    Show that 
    \bse 
        p^{\dot{\a}\a} := \widetilde{\l}^{\dot{\a}} \l^{\a} = (\bar{\sig})^{\dot{\a}\a}_{\mu} p^{\mu} = \begin{pmatrix}
            p^0 + p^3 & p^1 - ip^2 \\
            p^1 + ip^2 & p^0 - p^3
        \end{pmatrix}.
    \ese 
    This gives confirmation that \Cref{eqn:BispinorForm} and \Cref{eqn:pAlphaDotAlpha} agree.
\ebox 

If the momentum is real, then of course \Cref{eqn:BispinorForm} puts constraints on our $\widetilde{\l}/\l$s, in particular $\widetilde{\l} = \l^*$, up to some real phase (i.e. a factor 2 etc). However if $p\in\C$ then the $\widetilde{\l}$ and $\l$ are independent, this trick will be important later. You might make the argument that this is non-physical, which is true, however we will forget about this for calculational purposes. Namely we want to think of the $\l$s independently and then \textit{at the end} we will impose a reality condition.\footnote{This is an example of where we have injected a bit of twistor theory to motivate our calculations. The complex momentum comes from the fact that we consider complexified Minkowski spacetime and the momentum is given by derivatives w.r.t. the coordinates, which are now complex.} 

Note that this decomposition is far from unique. For example if we took 
\bse 
    \l \to e^{i\phi}\l \qand \widetilde{\l}\to e^{-i\phi} \widetilde{\l} 
\ese 
then $p$ is unchanged. This rescalling freedom is generated by a $U(1)$ generator, known as \textit{helicity}.\footnote{Its an example of a \textit{little group}.} Explicitly we have 
\be 
\label{eqn:HelicityOperator}
    h = -\frac{1}{2}\bigg(-\l^{\a}\frac{\p}{\p \l^{\a}} +\widetilde{\l}^{\dot{\a}}\frac{\p}{\p\widetilde{\l}^{\dot{\a}}}\bigg),
\ee 
which obeys 
\mybox{
    \be 
    \label{eqn:HelicityOfLambdas}
        h\l = -\frac{1}{2}\l \qand h\widetilde{\l} = \frac{1}{2}\widetilde{\l},
    \ee 
}
\noindent so undotted spinnors have helicity $-1/2$ and the dotted spinors have helicity $+1/2$. This basically describes a rotation in the plane orthogonal to the propagation direction; the $U(1)$ group generated by helicity corresponds to rotation in the 2-plane orthogonal to the \textit{spatial} momentum.\footnote{Note that it is $U(1)$ because we're working in 4d. That is the plane orthogonal to the spatial momentum is a 2-plane and so we just have a single angle and so it is abelian. In higher dimensions we have a 3-plane etc.}

We can raise/lower the indices using the Levi-Civita tensor 
\bse 
    \l_{\a} = \epsilon_{\a\beta}\l^{\beta} \qand \widetilde{\l}_{\dot{\a}} = \epsilon_{\dot{\a}\dot{\beta}}\widetilde{\l}^{\dot{\beta}}
\ese 
with
\be 
\label{eqn:LeviCivitaTensor}
    \epsilon_{\a\beta} = \epsilon_{\dot{\a}\dot{\beta}} = \begin{pmatrix}
        0 & -1 \\
        1 & 0
    \end{pmatrix}
\ee 
Note when you invert it you get a minus sign, namely 
\bse 
    \epsilon^{\a\beta} = \epsilon^{\dot{\a}\dot{\beta}} = \begin{pmatrix}
        0 & 1 \\
        -1 & 0
    \end{pmatrix}
\ese 
so that we have 
\bse 
    \epsilon_{\a\beta}\epsilon^{\beta\g} = \del^{\g}_{\a}. 
\ese 
and similarly for the dotted indices. From this we have the trivial self consistency check 
\bse 
    \l^{\a} = \epsilon^{\a\beta}\l_{\beta} = \epsilon^{\a\beta} \epsilon_{\beta\g} \l^{\g} = \del^{\a}_{\g}\l^{\g} = \l^{\a}.
\ese 

\br 
    For complete clarity, it only makes sense to contract an undotted index with an undotted one, and similarly dotted with dotted. That is we shouldn't be tricked by the equal signs in \Cref{eqn:LeviCivitaTensor} and multiply $\epsilon_{\a\beta}$ by $\epsilon^{\dot{\beta}\dot{\g}}$ and try contract the $\beta$ with $\dot{\beta}$ to get the identity matrix. In other words, the equal sign between $\epsilon_{\a\beta}$ and $\epsilon_{\dot{\a}\dot{\beta}}$ is mathematical none-sense: the two objects live in different spaces (the two $SU(2)$s) and so there is no notion of equality between them. This is probably a rather pedantic remark, but it is important to be clear on. 
\er

The idea of raising indices naturally allows us to define inner products on our two spaces.\footnote{In other words the Levi-Civita tensors are essentially giving us a map to the dual space, which can be used to define an inner product.} With the above remark in mind, it's important not to confuse the two separate inner products --- one for the undotted $SU(2)$ and one for the dotted $SU(2)$ --- and so we adopt the standard notation
\mybox{
    \be 
    \label{eqn:InnerProdcuts}
        \la i j \ra := \l_i^{\a}\l_{j\a} \qand [ij] := \widetilde{\l}_{i\dot{\a}}\widetilde{\l}_j^{\dot{\a}}, 
    \ee 
}
\noindent where the $i,j$ label which particle we're talking about (i.e. they label the momenta $p_i$). 

Recalling \Cref{rem:CommutingSpinors}, i.e. we are working with commuting spinors, we see that our inner products are antisymmetric:
\bse 
    \la i j \ra = - \la j i \ra \qand [ij] = - [ji]. 
\ese 
In particular it tells us that the inner product of a spinor with itself vanishes $\la ii \ra = 0 = [ ii ]$. This fact will be of great use to us in what follows and will allow us to drop a lot of terms. 

We have 
\bse 
	p_{\a\dot{\a}} = \epsilon_{\a\beta}\epsilon_{\dot{\a}\dot{\beta}}  p^{\dot{\beta}\beta} =  \l_{\a} \widetilde{\l}_{\dot{\a}}
\ese 
However we also have 
\bse 
    \epsilon_{\a\beta}\epsilon_{\dot{\a}\dot{\beta}}  p^{\dot{\beta}\beta} = \epsilon_{\a\beta}\epsilon_{\dot{\a}\dot{\beta}}  p_{\mu}(\bar{\sig}^{\mu})^{\dot{\a}\a} = p_{\mu} (\sig^{\mu})_{\a\dot{\a}} 
\ese 
where
\bse 
    (\sig^{\mu})_{\a\dot{\a}} = \epsilon_{\a\beta}\epsilon_{\dot{\a}\dot{\beta}} (\bar{\sig}^{\mu})^{\dot{\beta}\beta} = (\b1, \vec{\sig}). 
\ese 
which just gives us \Cref{eqn:pAlphaDotAlpha}.

\br 
    Note that when we use the Levi-Civita tensor to lower a combination of $\dot{\a}\a$ the order switches. That is we have $p^{\dot{\a}\a}$ with the dotted first but $p_{\a\dot{\a}}$ with the undotted first, and similarly for the $(\bar{\sig}^{\mu})^{\dot{\a}\a}$ and $(\sig^{\mu})_{\a\dot{\a}}$. 
\er 

Now consider the following
\bse 
    (p_i)_{\a\dot{\a}} (p_j)^{\dot{\a}\a} = \l_{i\a} \widetilde{\l}_{i\dot{\a}} \widetilde{\l}_j^{\dot{\a}}\l_j^{\a} = \la j i \ra [ij] = - \la i j \ra [ij],
\ese 
where again the $i,j$ label which particle we are talking about (i.e. they are not spatial components of the momenta). Now we can also write this in terms of the Pauli matrices 
\bse 
    \begin{split}
        p_i^{\mu}p_j^{\nu} (\sig_{\mu})_{\a\dot{\a}} (\bar{\sig}_{\nu})^{\dot{\a}\a} & = p_i^{\mu}p_j^{\nu}  \Tr[\sig_{\mu}\bar{\sig}_{\nu}] \\
        & = p_i^{\mu}p_j^{\nu} 2 \eta_{\mu\nu} \\
        & = 2p_i\cdot p_j \\
        & = (p_i+p_j)^2,
    \end{split}
\ese 
where we have used 
\bse 
    \Tr[\sig_{\mu}\bar{\sig}_{\nu}] = 2\eta_{\mu\nu},
\ese 
and where the last line follows from the fact that our momenta are null, $p_i^2=0=p_j^2$. Equating these two quantities gives us a very important formula:
\mybox{
    \be 
    \label{eqn:MomentumToInnerProducts}
        (p_i+p_j)^2 = 2p_i\cdot p_j = - \la ij \ra [ij]. 
    \ee 
}

To recap: essentially what \Cref{eqn:MomentumToInnerProducts} allows us to do is replace the kinematics (i.e. the momenta) of our null particles with inner products of spinor indices. This takes us part way to our stated at the beginning of this chapter. It is only part way as we still need to express the polarisations in terms of our spinors.

\section{Polarisations}

Polarisations of external particles are obtained from solutions to the classical equation of motion in the free theory. Let's start by considering the massless free Dirac Lagrangian\footnote{This is because the scalar fields don't have polarisations so the first starting point is spin-$1/2$.}
\bse 
    \cL = i \bar{\psi} \g^{\mu}\p_{\mu}\psi \qquad \bar{\psi} := \psi^{\dagger}\g^0,
\ese 
with 
\be 
\label{eqn:GammaMatrices}
    \g^{\mu} = \begin{pmatrix}
        0 & \sig^{\mu} \\
        \bar{\sig}^{\mu} & 0 
    \end{pmatrix}.
\ee 
The equations of motion of course the Dirac equation
\bse 
    i\g^{\mu} \p_{\mu} \psi = 0.
\ese 
Now consider a solution of the plane-wave form, namely 
\bse 
    \psi(x) = \begin{pmatrix}
        \chi_{\a}(p) \\
        \widetilde{\chi}^{\dot{\a}}(p)
    \end{pmatrix}  e^{ipx}.
\ese 
Plugging this into the Dirac equation gives us 
\bse 
    \g^{\mu}p_{\mu}\begin{pmatrix}
        \chi_{\a}(p) \\
        \widetilde{\chi}^{\dot{\a}}(p)
    \end{pmatrix} = 0 
\ese 
which with \Cref{eqn:GammaMatrices} reduces to
\bse 
    p_{\mu}\sig^{\mu}\widetilde{\chi} = 0 \qand p_{\mu}\bar{\sig}^{\mu} \chi = 0.
\ese

We now employ the work from the previous section and  write this in terms of spinors via
\bse 
    p_{\mu}(\sig^{\mu})_{\a\dot{\a}} = \l_{\a}\widetilde{\l}_{\dot{\a}} \qand p_{\mu} (\bar{\sig})^{\dot{\a}\a} = \widetilde{\l}^{\dot{\a}} \l^{\a}.
\ese 
From this, and the fact that our inner products are antisymmetric, we see that 
\bse 
    \chi_{\a} = \l_{\a} \qand \widetilde{\chi}^{\dot{\a}} = \widetilde{\l}^{\dot{\a}}
\ese 
is a valid solution (i.e. because $\l^{\a}\l_{\a}= 0 = \widetilde{\l}_{\dot{\a}}\widetilde{\l}^{\dot{\a}}$), then by standard uniqueness theorems we know this is the only solution. This motivates the following definitions. 
\mybox{
    \be 
    \label{eqn:Kets}
        \ket{p} := \begin{pmatrix}
            \l_{\a} \\
            0
        \end{pmatrix} \qand |p] := \begin{pmatrix}
            0 \\ 
            \widetilde{\l}^{\dot{\a}}
        \end{pmatrix}
    \ee 
    and 
    \be
    \label{eqn:Bras}
        \bra{p} := \begin{pmatrix}
            \l^{\a} & 0
        \end{pmatrix} \qand [p| := \begin{pmatrix}
            0 & \widetilde{\l}_{\dot{\a}}
        \end{pmatrix}
    \ee 
}
\noindent where the second set solve the equations of motion for $\bar{\psi}$, namely $\bar{\psi}\g^{\mu}p_{\mu}=0$. Note the placement of the $\a/\dot{\a}$ indices in \Cref{eqn:Kets,eqn:Bras}.

We take all the particles to be outgoing and make the following definitions 

\begin{center}
	\begin{tabular}{@{} C{4cm} C{4cm} C{4cm} @{}}
		\toprule
		Helicity & $+1/2$ & $-1/2$ \\ \\
		Quark & $[p|$ & $\bra{p}$ \\ \\
		Anti-Quark & $|p]$ & $\ket{p}$ \\
		\bottomrule
	\end{tabular}
\end{center}
so that the solution $[p|$ describes an outgoing quark with helicity $+1/2$, etc. If the momentum is physically \textit{ingoing} then we just make the substitution $p_{\mu} \to -p_{\mu}$ and $h\to - h$.

So we have a way to categorise the solutions to the free equations of motion in terms of the spinors. We therefore can express the polarisations in terms of the spinors which is exactly the goal we wanted to achieve. This was for spin-$1/2$ particles, i.e. Fermions, but the main focus of this course will be on spin-$1$ particles, Bosons, which we now discuss. 

\section{Spin-1}

As we just said, spin-$1$ particles correspond to Bosons, and in the standard model these in turn correspond to our propagators, e.g. the photon and gluons. Our main focus will be on gluons (and therefore QCD), with comments to other particles made.

\subsection{Abelian (QED)}

As always it is instructive to start off by discussion the abelian case, namely electrodynamics. We have the classical Maxwell action
\bse
    \cL = -\frac{1}{4}F_{\mu\nu}F^{\mu\nu} \qquad F_{\mu\nu} := \p_{\mu}A_{\nu} - \p_{\nu}A_{\mu}. 
\ese 
As we know from previous courses, this admits a gauge symmetry 
\bse 
    A_{\mu} \to A_{\mu} + \p_{\mu}\Lambda
\ese 
which can easily be used to show that $F_{\mu\nu}$ itself is invariant. We can use this freedom to impose \textit{Lorenz gauge}
\be
\label{eqn:LorenzGauge}
    \p_{\mu}A^{\mu} = 0.
\ee
In this guage, the equations of motion simplify to
\bse 
    0 = \p_{\mu}F^{\mu\nu} = \p_{\mu}\big(\p^{\mu}A^{\nu} - \p^{\nu}A^{\mu}\big) = \p^2 A^{\nu},
\ese 
which is just the wave equation 
\bse 
    \big(\p_t^2 - \nabla^2)A^{\nu} = 0.
\ese 
We now go to momentum space and again assume that $A_{\mu}(x)$ has a single plane wave solution 
\bse 
    A_{\mu}(x) = \epsilon_{\mu}(p) e^{ipx}. 
\ese 
Then our conditions above imply 
\bse 
    p^2 = 0 \qand \epsilon\cdot p = 0.
\ese 
where the second one comes from \Cref{eqn:LorenzGauge}. There is still some residual gauge freedom, however, notably 
\bse 
    A_{\mu} \to A_{\mu} + \p_{\mu}\omega \qquad \text{provided} \qquad \p^2 \omega = 0. 
\ese 
In momentum space this is 
\bse 
    \omega(x) = \a(p) e^{ipx} \qquad \implies \qquad p^2=0,
\ese 
then plugging this in above, we have 
\bse 
    \epsilon_{\mu}(p) \to \epsilon_{\mu}(p) + i \a(p) p_{\mu}
\ese
In other words $\epsilon$ and $\epsilon + i\a p$ are \textit{physically} equivalent. 

Let's choose a frame where the momentum points in the $z$ direction. 
\bse 
    p^{\mu} = (E,0,0,E)
\ese 
from which it follows that
\bse 
    \epsilon\cdot p = 0 \qquad \implies \qquad \epsilon^{\mu}=(0,b,c,0) + a(1,0,0,1).
\ese 
Then using our residual gauge freedom we can use
\bse
    \epsilon^{\mu} \to \epsilon^{\mu} - \frac{a}{E}p^{\mu} = (0,b,c,0).
\ese 
Finally imposing normalisation etc, we get the following basis for the polarisations
\mybox{
    \be
    \label{eqn:EpsilonPlusMinus}
        \epsilon_{\pm}^{\mu} = \frac{1}{\sqrt{2}} (0,1,\pm i, 0).
    \ee
}

Let's make some comments. The polarisation vectors are:
\ben[label=(\roman*)] 
    \item Complex conjugate related: $\epsilon_+= (\epsilon_-)^*$,
    \item Null: $\epsilon_+\cdot\epsilon_+ = 0 = \epsilon_-\cdot\epsilon_-$,
    \item Their inner product is $\epsilon_+\cdot \epsilon_- = -1$.
\een 
The key point is the second one because it tells us that our polarisation vectors are null, and so we can write them in terms of our spinors. We define
\mybox{
    \be 
    \label{eqn:PolarisationsAsSpinors}
        \begin{split}
            \epsilon_+^{\dot{\a}\a} := \epsilon_+^{\mu}(\bar{\sig}_{\mu})^{\dot{\a}\a} = -\sqrt{2} \frac{\widetilde{\l}^{\dot{\a}}\mu^{\a} }{\la \l\mu \ra } \\
            \epsilon_-^{\dot{\a}\a} := \epsilon_-^{\mu}(\bar{\sig}_{\mu})^{\dot{\a}\a} = \sqrt{2} \frac{\widetilde{\mu}^{\dot{\a}}\l^{\a}} {[\widetilde{\l}\widetilde{\mu}] }
        \end{split}
    \ee 
}
\noindent where $p^{\dot{\a}\a} = \widetilde{\l}^{\dot{\a}} \l^{\a}$ and $q^{\dot{\a}\a} = \widetilde{\mu}^{\dot{\a}}\mu^{\a}$ is a reference momentum that encodes the residual gauge symmetry. We will prove this in just a moment. First let's check that these give us the desired properties. We have 
\bse 
    (\epsilon_+)^{\dot{\a}\a} (\epsilon_+)_{\a\dot{\a}} \sim [\widetilde{\l}\widetilde{\l}] = 0 \qand (\epsilon_-)^{\dot{\a}\a} (\epsilon_-)_{\a\dot{\a}}  \sim \la \l \l \ra = 0,
\ese
which is condition (ii). 

Next we have
\bse 
    (\epsilon_+)^{\dot{\a}\a}(\epsilon_-)_{\a\dot{\a}} = -2 \frac{[\widetilde{\mu}\widetilde{\l}]\la \mu\l\ra }{\la \l \mu \ra [\widetilde{\l}\widetilde{\mu}]} = -2.
\ese 
which seems wrong at first (we want $-1$ not $-2$) but then we remember that the $-1$ condition comes from the Lorentz expressions. As we see from \Cref{eqn:MomentumToInnerProducts} the two are related by a factor of $2$, i.e.
\be
\label{eqn:LorentzToSpinorContraction}
    2 \epsilon_+\cdot\epsilon_- = (\epsilon_+)^{\dot{\a}\a}(\epsilon_-)_{\a\dot{\a}},
\ee 
and so we get the right-hand side inner product being $-1$. 

We also want to check that our gauge condition $\epsilon\cdot p =0$ is preserved. This follows easily from 
\bse
    \epsilon_+^{\dot{\a}\a} p_{\dot{\a}\a} \sim [\widetilde{\l}\widetilde{\l}] = 0 \qand \epsilon_-^{\dot{\a}\a} p_{\dot{\a}\a} \sim \la \l \l \ra = 0.
\ese 
Equally we can check the helicity, using \Cref{eqn:HelicityOperator} we can easily obtain
\bse 
    h \epsilon_{\pm} = \pm \epsilon_{\pm},
\ese 
which is where the subscript came from in the first place. 

Finally we want to show that $\mu$ does indeed encode the residual gauge symmetry. We consider a variation of $\mu$ and compute how $\epsilon$s change. 
\bse 
    \mu \to \mu + a \mu + b \l  = (1+a)\mu + b\l,
\ese 
where we have used that we can use $\mu$ and $\l$ as a basis of two component objects, which is fine as we have already assumed $\mu$ and $\l$ are not proportional. So how does $\epsilon$ change? We have (using $\la \l \l \ra = 0$)
\bse
    \begin{split}
        (\epsilon_+)^{\dot{\a}\a} &\to -\sqrt{2} \frac{ \widetilde{\l}^{\dot{\a}} \big((1+a)\mu^{\a} + b\l^{\a}\big) }{\la \l, (1+a)\mu + b\l \ra} \\
        & = -\sqrt{2} \frac{\widetilde{\l}^{\dot{\a}}\mu^{\a} }{\la \l \mu\ra} - \frac{\sqrt{2}b}{1+a}\frac{\widetilde{\l}^{\dot{\a}} \l^{\a} }{\la \l \mu\ra} \\
        & = (\epsilon_+)^{\dot{\a}\a} - \frac{\sqrt{2}b}{(1+a)\la \l \mu \ra} p^{\dot{\a}\a},
    \end{split}
\ese 
which is just a residual gauge transformation, i.e. $(\epsilon_+)^{\dot{\a}\a}$ has changed by something proportional to $p$. A similar calculation will give the $(\epsilon_-)^{\dot{\a}\a}$ result.


\chapter{Colour-Ordering}

As we made clear before starting the calculation, the above results are for \textit{abelian} gauge theories, i.e. QED. We obviously now want to ask the question of how this translates to non-abelian theories in order to study QCD. Indeed QCD will be the main focus of this course, for the simple reason that we know QCD amplitudes are tedious to calculate from the Feynman diagram approach (mainly due to the gluon self interactions). 

The polarisation story from the QED case before will all follow through, so all we really need to consider here is the colour part of it. The idea is to split the amplitude into a colour part and a kinematical part and then just put them together in the end. 

\section{Setting Up The Lagrangian}

So how do we do this? Well as always we `promote' our gauge fields $A_{\mu}$ to be matrices 
\bse 
    {(A_{\mu})_i}^j = A_{\mu}^a {(T^a)_i}^j
\ese
where $T^a$ are the generators of $SU(N)$ and $a \in \{1,2,...,N^2-1\}$ is an adjoint index. We shall use the convention\footnote{Note this is different to the other standard convention $\Tr[T^aT^b]=\frac{1}{2}\del^{ab}$. For $SU(2)$ the one used here corresponds to using $T^a = \frac{1}{\sqrt{2}}\sig^a$, where $\sig^a$ are the Pauli matrices, while the other convention uses $\tau^a = \frac{1}{2}\sig^a$. Similarly for other $SU(N)$s with the Pauli matrices respectively replaced.}
\be
\label{eqn:TraceTaTb}
    \Tr[T^aT^b] = \del^{ab}.
\ee 
This allows us to extract the individual $A_{\mu}^a$s simply via 
\be 
\label{eqn:AmuaFromAmu}
    A_{\mu}^a = \Tr[A_{\mu}T^a].
\ee 
The commutator of our generators obey 
\bse 
    [T^a, T^b] = i\sqrt{2} f^{abc}T^c,
\ese
where the perhaps unfamiliar $\sqrt{2}$ factor comes from using convention \Cref{eqn:TraceTaTb}.

Our Lagrangian now simply becomes\footnote{Note this is actually the same as for the abelian case; in that case there is only one term in the trace so we don't need to write it.}
\bse 
    \cL = -\frac{1}{4}\Tr[ F_{\mu\nu}F^{\mu\nu}]
\ese 
with 
\bse 
    F_{\mu\nu} = \p_{\mu}A_{\nu} - \p_{\nu}A_{\mu} - \frac{ig}{\sqrt{2}}[A_{\mu},A_{\nu}] 
\ese
where we explicitly see the non-ableian term, the commutator. We can write this in terms of $F_{\mu\nu}^a$ using the structure constants
\bse 
    F_{\mu\nu}^a = \p_{\mu}A_{\nu}^a - \p_{\nu}A_{\mu}^a + gf^{abc}A^b_{\mu}A^c_{\nu}. 
\ese 
If we introduce the covariant derivative 
\bse 
    D_{\mu} = \p_{\mu} - \frac{ig}{\sqrt{2}} A_{\mu} 
\ese 
we can write the field strength simply as 
\bse 
    F_{\mu\nu} = \frac{\sqrt{2}i}{g} [D_{\mu},D_{\nu}]
\ese 

We know that the Lagrangian has the gauge symmetry 
\bse 
    A_{\mu} \to U A_{\mu} U^{\dagger} + \frac{i}{g}U(\p_{\mu}U^{\dagger}) \qquad U \in SU(N).
\ese
This gives the transformation 
\bse 
    F_{\mu\nu} \to U F_{\mu\nu} U^{\dagger}
\ese
which is easily seen from checking that\footnote{Really we shouldn't say that $D_{\mu}$ transforms in this way but rather that $D_{\mu}\psi \to U D_{\mu} \psi$ where $\psi$ is a field that transforms as $\psi \to U \psi$. However it is standard to write it like this and the idea is clear.}
\bse 
    D_{\mu} \to U D_{\mu} U^{\dagger}.
\ese 

\br 
    Note that we were careful to say that the \textit{Lagrangian} is gauge invariant, not the field strength itself. In other words the field strength is gauge \textit{covariant}. This is easily checked using the cyclicity of the trace and $U^{-1}=U^{\dagger}$ for $SU(N)$. This point is just highlighted because in the abelian case we had that the field strength itself was gauge invariant.
\er 

Next recall that this gauge transformation introduces a redundancy into the system. We fix this redundancy in the path integral formulation by using the Fadeev-Popov procedure.\footnote{See the QFT II course for more details.} This essentially boils down to adding a new term to the action 
\bse 
    \cL = - \frac{1}{4}\Tr[F_{\mu\nu}F^{\mu\nu}] - \frac{1}{2}\Tr[G G] + \cL_{\text{ghosts}},
\ese 
where $G(x)$ is a \textit{gauge fixing function}. The ghost terms we will neglect as we are only concerned with amplitudes (which only know about external on-shell particles) and so the ghosts never appear (they only flow through loops). 

For the gauge fixing condition we will choose 
\mybox{
\be 
\label{eqn:GervaisNeveuGauge}
    G(x) = \p_{\mu} A^{\mu} - \frac{ig}{\sqrt{2}} A_{\mu}A^{\mu}
\ee 
}
\noindent this is known as the \textit{Gervais-Neveu gauge}. 

\br 
    Note that if we set $g=0$ then we just get a Lorenz gauge result. 
\er 

The reason we take this choice is because then the Lagrangian simply becomes 
\bse 
    \cL = \Tr[ \frac{1}{2} A_{\mu}\p^2 A^{\mu} - i \sqrt{2} g \p^{\mu}A^{\nu}A_{\nu} A_{\mu} + \frac{1}{4}g^2 A^{\mu}A^{\nu}A_{\mu}A_{\nu}]. 
\ese 

\section{Feynman Diagrams In Double Line Notation}

We want to get the Feynman rules for this. We shall use a "quick and dirty" way to get them. We will obtain the Feynman rules in terms of \textit{matrix valued} expressions. This is different to what we normally do, where we normally have explicit colour indices. In other words, normally we obtain Feynman rules for the $A_{\mu}^a$s, but here we are going to write them down for ${(A_{\mu})_i}^j = A_{\mu}^a (T^a{)_i}^j$. We will of course need to keep track of how the matrix indices are contracted and this will give rise to what are known as \textit{ribbon diagrams}, which use so-called \textit{double line notation}. 

We start with the propagator. This is extracted from the kinetic term:\footnote{Our gluons are massless so we only need this term. That is normally we also consider the mass term when finding the propagator.}
\bse 
    \frac{1}{2}{(A_{\mu})_i}^j \p^2 {(A^{\mu})_j}^i = \frac{1}{2} {(A_{\mu})_i}^j \p^2 {(A_{\nu})_m}^{\ell} \eta^{\mu\nu} \del^i_{\ell} \del^m_j,
\ese 
where the right hand side just shows the contractions explicitly. We depict this in terms of the double line diagrams:
\begin{center}
    \btik 
        \midarrow (1,0.2) -- (-1,0.2);
        \midarrow (-1,-0.2) -- (1,-0.2) node [midway, below] {$p$};
        \node at (-1.2,-0.2) {$j$};
        \node at (1.2,0.2) {$\ell$};
        \node at (-1.2,0.2) {$i$};
        \node at (1.2,-0.2) {$m$};
        \node at (-1.5,0) {$\mu$};
        \node at (1.5,0) {$\nu$};
        \node[right] at (2,0) {\Large{$=$}};
        \node[right] at (3,0.1) {\huge{$- \frac{ i \eta^{\mu\nu} \del_{\ell}^i \del_j^m }{p^2+i\epsilon}$}};
    \etik 
\end{center}

Let's clear up any potential confusion on how to construct such a diagram and obtain the mathematical expression from it.
\begin{itemize}
    \item This is not two particles propagating but a single propagator. The two lines correspond to the fact we have $2$ $A$s.
    \item You label the end of each double line with the Lorentz structure, $\mu/\nu$ 
    \item You label each the end of each line with the matrix indices, $i,j$ etc.
    \item You join contracted indices with the convention that the arrow points from the lower index on the $\del$ to the upper index on the $\del$. This corresponds to pointing from the upper index on one $A$ to the lower index on the other $A$. Note this will always result in adjacent lines pointing in opposite directions.
    \item You then obtain the mathematical expression from the above with the usual Feynman rules:
        \ben[label=(\alph*)]
            \item Factor of $i$ from $e^{iS}$
            \item Minus sign from $\p^2 \to -p^2$ 
            \item Contour argument to get $+i\epsilon$ in denominator
        \een 
\end{itemize}

Next we look at the three point vertex. We now have to be a bit more careful because more things are contracted and ordering matters. First let's write down the term in the Lagrangian that gives rise to it in matrix form:
\bse 
    -i\sqrt{2} \, g \, \p_{\mu} (A_{\nu}{)_k}^j (A^{\nu}{)_j}^i (A^{\mu}{)_i}^k
\ese 
Now follow the proscription given above to draw:

\begin{center}
    \btik[scale=1.5]
        \midarrow (-1,-1) -- (-0.2,0);
        \midarrow (-0.2,0) -- (-0.2,1);
        \midarrow  (0.2,0) -- (1,-1);
        \midarrow (0.2,1) -- (0.2,0);
        \midarrow (0.6,-1) -- (0,-0.2);
        \midarrow (0,-0.2) -- (-0.6,-1);
        % 
        \node at (-1.1,-1.1) {$n$};
        \node at (-0.7,-1.1) {$m$};
        \node at (1.1,-1.1) {$k$};
        \node at (0.7,-1.1) {$\ell$};
        \node at (0.2,1.2) {$j$};
        \node at (-0.2,1.2) {$i$};
        %
        \node at (-0.9,-1.4) {\large{$\mu,p$}};
        \node at (0.9,-1.4) {\large{$\rho,r$}};
        \node at (0,1.5) {\large{$\nu,q$}};
    \etik   
\end{center}
where we have now labelled the momentum next to the Lorentz index. Mathematically this diagram corresponds to 
\bse 
    i (-i\sqrt{2}g) \del^i_n \del^m_{\ell} \del^k_j \big( (-iq_{\nu})  \eta_{\nu\rho} + \text{cyclic permutations} \big) = -i \sqrt{2} \, g \, (q_{\mu}\eta_{\nu\rho} + r_{\nu}\eta_{\rho\mu} + p_{\rho} \eta_{\mu\nu}) \del_n^i \del_j^k \del_{\ell}^m,
\ese 
where we get the momentum terms from the derivative term. 

Finally let's look at the four point interaction:
\bse 
    \frac{1}{4}g^2 (A^{\mu}{)_j}^i (A^{\nu}{)_i}^k (A_{\mu}{)_k}^{\ell}  (A_{\nu}{)_{\ell}}^j 
\ese 
This is a little nicer because it doesn't contain a derivative. The diagram is simply 
\begin{center}
    \btik
        \midarrow (-1.7,-1.5) -- (-0.2,0);
        \midarrow (-0.2,0) -- (-1.7,1.5);
        \node at (-1.9,-1.6) {$q$};
        \node at (-1.9,1.6) {$i$};
        \begin{scope}[rotate around={90:(0,0)}]
            \midarrow (-1.7,-1.5) -- (-0.2,0);
            \midarrow (-0.2,0) -- (-1.7,1.5);
            \node at (-1.9,-1.6) {$n$};
            \node at (-1.9,1.6) {$p$};
        \end{scope}
        \begin{scope}[rotate around={-90:(0,0)}]
            \midarrow (-1.7,-1.5) -- (-0.2,0);
            \midarrow (-0.2,0) -- (-1.7,1.5);
            \node at (-1.9,-1.6) {$j$};
            \node at (-1.9,1.6) {$k$};
        \end{scope}
        \begin{scope}[rotate around={180:(0,0)}]
            \midarrow (-1.7,-1.5) -- (-0.2,0);
            \midarrow (-0.2,0) -- (-1.7,1.5);
            \node at (-1.9,-1.6) {$\ell$};
            \node at (-1.9,1.6) {$m$};
        \end{scope}
        % 
        \node at (-2.1,2) {\large{$\mu$}};
        \node at (-2.1,-2) {\large{$\l$}};
        \node at (2.1,2) {\large{$\nu$}};
        \node at (2.1,-2) {\large{$\rho$}};
        \node[right] at (2.5,0) {\Large{$= \, i g^2 \del^i_q \del^k_j \del^m_{\ell} \del^p_n \eta_{\mu\rho}\eta_{\nu\l} $}};
    \etik
\end{center}
Note that there is no factor of $4$ from the cyclic permutations. 

\subsection{Colour-Ordered Amplitudes}

We can now compute some amplitudes, and we start with the 3-point amplitude. The idea is to insert a factor of $T^a$ at each vertex and take the trace, this then allows us to extract the field $A_{\mu}^a$ via \Cref{eqn:AmuaFromAmu}. We then also have to include the polarisation vector, so in total we have:
\begin{center}
    \btik 
        \draw[thick, ->] (0,0) -- (0,1.5) node [above] {$a_1,p_1,\epsilon_1$};
        \draw[thick, ->, rotate around={120:(0,0)}] (0,0) -- (0,1.5) node [below] {$a_3,p_3,\epsilon_3$};
        \draw[thick, ->, rotate around={-120:(0,0)}] (0,0) -- (0,1.5) node [below] {$a_2,p_2,\epsilon_2$};
        \node at (3.5,0.5) {\Large{$=$}};
        \begin{scope}[xshift=7cm, yshift=0.25cm]
            \midarrow (-1,-1) -- (-0.2,0);
            \midarrow (-0.2,0) -- (-0.2,1);
            \midarrow  (0.2,0) -- (1,-1);
            \midarrow (0.2,1) -- (0.2,0);
            \midarrow (0.6,-1) -- (0,-0.2);
            \midarrow (0,-0.2) -- (-0.6,-1);
            % 
            \node at (-1.1,-1.1) {$n$};
            \node at (-0.7,-1.1) {$m$};
            \node at (1.1,-1.1) {$k$};
            \node at (0.7,-1.1) {$\ell$};
            \node at (0.2,1.2) {$j$};
            \node at (-0.2,1.2) {$i$};
            %
            \node at (-0.9,-1.4) {$\mu$};
            \node at (0.9,-1.4) {$\rho$};
            \node at (0,1.5) {$\nu$};
            %
            \node at (0,2) {$(T^{a_1}{)_i}^j \epsilon_1^{\nu}$};
            \node at (1.2,-1.8) {$(T^{a_2}{)_k}^{\ell} \epsilon_2^{\rho}$};
            \node at (-1.2,-1.8) {$(T^{a_3}{)_m}^n \epsilon_3^{\mu}$};
            \node[right] at (2,0) {\large{$+$ perms of $(1,2,3)$}};
        \end{scope}
    \etik 
\end{center}

We can then use this to obtain 
\bse 
    iT_3 = -i\sqrt{2}g \Tr[T^{a_1}T^{a_2}T^{a_3}] \big( \epsilon_1\cdot \epsilon_2 p_1\cdot \epsilon_3 + \epsilon_2\cdot\epsilon_3 p_2\cdot \epsilon_1 + \epsilon_3\cdot \epsilon_1 p_3 \cdot \epsilon_2\big) + (2 \leftrightarrow 3),
\ese 
where the second part is to account for the non-cyclic permutations; the cyclic permutations are already taken care of in our Feynman rules. 

We did this calculation for the 3-point vertex, but the claim is that more generally, using the double line notation, we see that all \textit{tree-level} amplitudes can be written as 
\mybox{
    \be
    \label{eqn:iTn}
        iT_n = g^{n-2} \sum_{\substack{\text{non-cyclic} \\ \text{perms}}} \Tr[T^{a_1} ... T^{a_n}] A(1,...,n)
    \ee 
}
\noindent where the $A(1,...,n)$ is a purely kinematical thing known as the \textit{colour ordered amplitude} and is sum over all diagrams with a fixed cyclic ordering of the gluons and no crossed lines. The colour ordered amplitude knows nothing about the group theoretic piece (i.e. the colour structure), which is all in the trace. The sum, which is over non-cyclic permutations, can be expressed as the sum over permutations of legs $2,...,n$ holding leg $1$ fixed. 

This is not a trivial result so let's summarise the main points: 
\ben[label=(\roman*)] 
    \item We can reduce the calculation of an $n$-point \textit{tree level} amplitude to the calculation of colour-ordered amplitudes. We then get the full amplitude by "\textit{colour-dressing}": multiply by a cyclically ordered colour factor (the trace) and sum over non-cyclic permutations.
    \item The colour ordered amplitudes are cyclically symmetric,
    \bse 
        A(1,2,...,n) = A(2,...,n,1),
    \ese 
    which can be seen from the fact that the trace is cyclically symmetric. 
    \item We compute the colour ordered amplitudes using our double line notation Feynman rules derived before. We state these colour ordered Feynman rules in the next table.
\een 

\mybox{
    \begin{center}
    	\begin{tabular}{@{} C{4cm} C{4cm} C{4cm} @{}}
    		\toprule
    		Name & Diagram & Math Expression \\ 
    		\midrule
    		Propagator & \btik 
    		    \midarrow (0,0) -- (1.5,0) node [midway, below] {$p$};
    		    \node at (-0.2,0) {$\mu$};
    		    \node at (1.7,0) {$\nu$};
    		\etik & \bse 
    		    -\frac{i\eta_{\mu\nu}}{p^2+i\epsilon}
    		\ese \\ 
    		3-Point Vertex & \btik 
    		    \midarrow (0,0) -- (0,1) node [above] {$\mu_1,p_1$};
    		    \midarrow[rotate around={120:(0,0)}] (0,0) -- (0,1) node [below] {$\mu_3,p_3$};
    		    \midarrow[rotate around={-120:(0,0)}] (0,0) -- (0,1) node [below] {$\mu_2,p_2$};
    		\etik & \bse
    		    \begin{split}
    		           -i \sqrt{2} \big( \eta^{\mu_1\mu_2} p_1^{\mu_3} \\
    		           \quad + \eta^{\mu_2\mu_3} p_2^{\mu_1} \\
    		           \quad + \eta^{\mu_3\mu_1} p_3^{\mu_2} \big)
    		    \end{split}
    		\ese \\
    		4-Point Vertex & \btik 
    		    \midarrow[rotate around={45:(0,0)}] (0,0) -- (0,1) node [above] {$\mu_1$};
    		    \midarrow[rotate around={-45:(0,0)}] (0,0) -- (0,1)  node [above] {$\mu_2$};
    		    \midarrow[rotate around={45:(0,0)}] (0,0) -- (0,-1)  node [below] {$\mu_3$};
    		    \midarrow[rotate around={-45:(0,0)}] (0,0) -- (0,-1)  node [below] {$\mu_4$};
    		\etik  & \bse 
    		    i \eta^{\mu_1\mu_3} \eta^{\mu_2\mu_4}
    		\ese \\
    		\bottomrule
    	\end{tabular}
    \end{center}
}

Note we have dropped the factors of $g$ and the $\del$s. We are no longer using double line notation, since colour ordered amplitudes do not contain information about colour. It is understood that these vertices will be used to construct diagrams with fixed cyclic ordering of external legs and no crossed lines. 
\chapter{$n$-Point Amplitudes}

So far we have derived two very powerful techniques: the spinor-helicity formalism and colour-ordering. We now want to put these two techniques together in order to compute some amplitudes.

\section{3-Point Amplitudes}

We start by considering $3$-point amplitudes as these are the most simple and give hints at how to extend to higher point amplitudes. The first thing we note is that for 3-point amplitudes \textit{all} kinematical invariants vanish. We see this simply from 
\bse 
    p_1+p_2+p_3 = 0 \qquad\implies \qquad 2p_1\cdot p_2 = (p_1+p_2)^2 = p_3^2 = 0,
\ese 
and similarly $p_1\cdot p_3 = 0 = p_2\cdot p_3$. If we convert this into spinor formalism, it becomes
\be 
\label{eqn:3PointSpinorInnerProductVanish}
    0 = (p_i+p_j)^2 = - \la ij \ra [ij].
\ee 
Next we recall that if we have real momenta we require $\la ij \ra = -[ij]^*$, and so the only way we can satisfy the above relation is if all spinor inner products vanish. It follows from this that we cannot define any non-trivial amplitudes. Ah... that's not so good.

So what do we do? Well the above followed directly from us imposing that the momenta are real, so we could ask the question "what if they're complex?" As we said before, in this $\l^{\a}$ and $\widetilde{\l}^{\dot{\a}}$ are independent and so the two inner products are independent. This means we could satisfy \Cref{eqn:3PointSpinorInnerProductVanish} by setting \textit{either} $\la i j\ra =0$ \textit{or} $[ij]=0$, with the other not needing to vanish. This is exactly what we are going to do. 

Before moving on it is worth keeping in mind that this complex momenta condition comes from complexifying our spacetime and so is unphysical. Therefore the results we obtain here do not represent good physics (they will all vanish when we impose reality conditions), however they are worth pursuing as they give nice mathematical insight. 

The first thing we do is to express our momentum conservation in terms of spinors. Using $p_i = \l_i \widetilde{\l}_i$ we have 
\bse 
    \l_1\widetilde{\l}_1 + \l_2\widetilde{\l}_2 + \l_3\widetilde{\l}_3 = 0.
\ese 
Next let's contract with $\l_1$, using $\la 11 \ra = 0$, we are left with
\bse 
    \la 12 \ra \widetilde{\l}_2 + \la 13 \ra \widetilde{\l}_3 = 0.
\ese 
If we then assume that $\la 12 \ra \neq 0$, we obtain 
\bse 
    \widetilde{\l}_2 = \frac{\la 13 \ra}{\la 12 \ra} \widetilde{\l}_3.
\ese 
Similarly if we had contracted with $\l_2$ we would obtain 
\bse 
    \widetilde{\l}_1 = \frac{\la 23 \ra }{\la 21 \ra }\widetilde{\l}_3,
\ese
and so in total we have 
\bse 
    \widetilde{\l}_1 \propto \widetilde{\l}_2 \propto \widetilde{\l}_3,
\ese 
but if this is the case then we have 
\bse 
    [ij] \propto [ii] = 0 
\ese 
and so all square bracket inner products vanish. That is 
\bse 
    [12] = [13] = [23] = 0.
\ese 
Similarly if we had taking contractions with the tilded $\widetilde{\l}$s and assumed that the square bracket inner products were non-vanishing we would obtain 
\bse 
    \la 12 \ra = \la 13 \ra = \la 23 \ra = 0.
\ese
So we have that all 3-point amplitudes \textit{only} depend on $\la ij \ra $ \textit{or} only depend on $[12 ]$. Note this result is more restrictive then \Cref{eqn:3PointSpinorInnerProductVanish} itself, as this is satisfied with
\bse 
    \la 12 \ra = [23] = [13] =0,
\ese
and so the 3-point amplitude would depend on $[12], \la 23 \ra$ and $\la 13\ra$. 

\br 
    The above result, that the 3-point amplitude only depends on $\la ij \ra$ or $[ij]$, turns out to be equivalent to saying that they are holomorphic or antiholomorphic, respectively. 
\er

Ok let's compute the 3-point amplitudes explicitly. From our colour ordered Feynman rules, we have 
\begin{center}
    \btik 
        \node[left] at (-1.5,0) {\large{$A(\{p_i,\epsilon_i\}) \, =$}};
        \midarrow (0,0) -- (0,1) node [above] {$\epsilon_1,p_1$};
        \midarrow[rotate around={120:(0,0)}] (0,0) -- (0,1) node [below] {$\epsilon_3,p_3$};
        \midarrow[rotate around={-120:(0,0)}] (0,0) -- (0,1) node [below] {$\epsilon_2,p_2$};
        \node[right] at (1.5,0) {\large{$= \, \sqrt{2}\big[ (\epsilon_1\cdot \epsilon_2)(\epsilon_3 \cdot p_1) + (\epsilon_2\cdot \epsilon_3)( \epsilon_1 \cdot p_2) + (\epsilon_3\cdot \epsilon_1)(\epsilon_2 \cdot p_3) \big]$}};
    \etik
\end{center}
where we have 
\bse 
    iT_3 = -ig^{3-2}A = -igA. 
\ese

We now want to simplify this. We do this by first considering the case where all the gluons have the same helicity. W.l.o.g. let's take them to all be $+$. Then from \Cref{eqn:PolarisationsAsSpinors} we have
\bse 
    \epsilon_+(p_i,q_i) \cdot \epsilon_+(p_j,q_j) = 2 \frac{\la q_i q_j\ra [p_jp_i]}{\la q_i p_i\ra \la q_j p_j \ra}
\ese 
where\footnote{Again here $i$ labels the gluon we're considering, not a spatial index.} $p_i^{\dot{\a}\a} = \widetilde{\l}^{\dot{\a}}_i\l_i^{\a}$ are the momenta and $q_i^{\dot{\a}\a} = \widetilde{\mu}^{\dot{\a}}_i\mu^{\a}_i$ are the reference momenta encoding the gauge invariance. The inner products in the above are meant to be understood as the relative parts, i.e. 
\bse 
    \la q_i q_j \ra =  \mu_i^{\a} \mu_{j\a} \qand [p_j p_i] = \widetilde{\l}_{j\dot{\a}} \widetilde{\l}^{\dot{\a}}_i,
\ese
etc. Now if we choose $q_1=q_2=q_3$ then $\la q_i q_j \ra = 0$ and so the above vanishes. A similar calculation shows that 
\bse 
    \epsilon_-(p_i,q_i) \cdot \epsilon_-(p_j,q_j) = 0.
\ese    
We summarise this below:
\mybox{
    \be 
    \label{eqn:Apmpmpm}
        A(\pm \pm \pm) =0,
    \ee 
}
\noindent where hopefully the notation is clear.

\subsection{MHV \& $\overline{\text{MHV}}$}

Ok so what if only one polarisation is different? That is consider $A(1^-2^-3^+)$, which is the 3-point minimal helicity violating (MHV) amplitude. Well from above we have 
\be 
\label{eqn:A--+}
    A(1^-,2^-,3^+) = \sqrt{2}\big[ (\epsilon_1^-\cdot \epsilon_2^-)(\epsilon_3^+ \cdot p_1) + (\epsilon_2^-\cdot \epsilon_3^+)( \epsilon_1^- \cdot p_2) + (\epsilon_3^+\cdot \epsilon_1^-)(\epsilon_2^- \cdot p_3) \big],
\ee
Next, we have\footnote{Note on the minus sign in the $\epsilon^+_i\cdot \epsilon^-_j$ expression is taken care of by swapping $[q_jp_i]=-[p_iq_j]$. Similarly for the $\epsilon^+_i\cdot p_j$ with $[p_jp_i]=-[p_ip_j]$.}
\be 
\label{eqn:PolarisationInnerProducts}
    \begin{split}
        \epsilon_i^- \cdot \epsilon_j^- & = \frac{[q_iq_j] \la p_jp_i\ra }{[q_ip_i] [q_j p_j]} \\
        \epsilon_i^+ \cdot \epsilon_j^- & = \frac{ [p_iq_j] \la q_ip_j\ra }{\la q_ip_i\ra [q_j p_j] \ra} \\
        \epsilon_i^+ \cdot p_j & = \frac{ [p_ip_j] \la q_ip_j\ra }{ \sqrt{2}\la q_ip_i\ra} \\
        \epsilon_i^- \cdot p_j & = \frac{ [p_jq_i] \la p_i p_j \ra }{\sqrt{2} [q_ip_i] }
    \end{split}
\ee 
where we have used the fact that when going from Lorentz contractions (which the $\cdot$ represents) to spinor contractions we get a factor of $1/2$, as in \Cref{eqn:LorentzToSpinorContraction}. We can then express our 3-point MHV amplitude in terms of spinor inner products. Again the $q_i^{\dot{\a}\a} = \widetilde{\mu}^{\dot{\a}} \mu^{\a}$ are our reference momenta. If we pick 
\bse 
    \widetilde{\mu}_1 = \widetilde{\mu}_2 \qand \mu_3 = \l_1
\ese 
then we have 
\bse 
    [q_1q_2] = 0 \qand \la q_3 p_1 \ra = 0.
\ese 
From this only the middle term in \Cref{eqn:A--+} survives:
\bse 
    A(1^-,2^-,3^+) = \frac{\la q_3 2 \ra [3q_2]}{\la q_3 3 \ra} \frac{ [q_1 2] \la 21 \ra }{[1q_1]},
\ese 
where we have used a notation where just a number means $p$, i.e. $\la q_3 2 \ra = \la q_3 p_2\ra$. Then we use our reference momenta definitions to replace 
\bse 
    \begin{split}
        \la q_3 2 \ra & \longrightarrow \la 12 \ra \\
        \la q_3 2 \ra & \longrightarrow \la 12 \ra \\
        [q_1 2] & \longrightarrow [q_1 2] \\
        [1 q_1] & \longrightarrow [1q_2],
    \end{split}
\ese 
to obtain 
\bse 
    A(1^-,2^-,3^+) = - \frac{\la 12\ra^2}{\la 13\ra } \frac{[3q_2]}{[1q_2]}
\ese 
where we have also used $\la 21 \ra = -\la 12 \ra$. 

This is nice, but really we want to remove the reference momenta from the expression. We do this by considering momentum conservation 
\bse 
    \l_1 \widetilde{\l}_1 + \l_2 \widetilde{\l}_2 + \l_3 \widetilde{\l}_3 = 0,
\ese 
which if we contract with $\l_2 \widetilde{\mu}_2$ we get 
\be 
\label{eqn:3q21q2}
    \la 21 \ra [1q_2] + \la 23 \ra [3 q_2] = 0 \qquad \implies \qquad \frac{[3q_2]}{[1q_2]} = \frac{\la 12\ra }{\la 23 \ra}. 
\ee 
Putting this all together we get 
\mybox{
    \be 
    \label{eqn:A--+Final}
        A(1^-,2^-,3^+) = \frac{\la 12 \ra^4}{\la 12 \ra \la 23 \ra \la 31 \ra}, \qquad \text{with} \qquad \widetilde{\l}_1 \propto \widetilde{\l}_2 \propto \widetilde{\l}_3,
    \ee 
}
\noindent where we have multiplied denominator and numerator by $\la 12\ra$ to get a more symmetric looking denominator. Similarly we can calculate the \textit{anti-MHV} (or $\overline{\text{MHV}}$) expression 
\mybox{
    \be 
    \label{eqn:A++-Final}
        A(1^+,2^+,3^-) = -\frac{ [12]^4}{[12][23][31]}, \qquad \text{with} \qquad \l_1 \propto \l_2 \propto \l_3.
    \ee 
}

Let's make a couple comments:
\begin{itemize}
    \item We now note that we only set $\widetilde{\mu}_1=\widetilde{\mu}_2$ but didn't say how they were related to the $\widetilde{\l}_i$s. Now if we set $\widetilde{\mu}_2 = \widetilde{\l}_3$ then we would have $[3q_2]=0$ and so the numerator in \Cref{eqn:3q21q2} would vanish, which in turn suggest that \Cref{eqn:A--+Final} vanishes. However we then note that we already have $\widetilde{\mu}_1 = \widetilde{\mu}_2$ and so we must also have $\widetilde{\mu}_1 = \widetilde{\l}_3 \propto \widetilde{\l}_1$, and so $[1q_2]=0$. This is the denominator of \Cref{eqn:3q21q2}, and so we just get $\frac{0}{0}$, which is ill-defined.  
    \item For real momenta we know all our inner products vanish and so our amplitudes simply become 
    \bse 
        A(1^-,2^-,3^+) = \frac{0^4}{0^3} = 0 \qand A(1^+,2^+,3^-) = - \frac{0^4}{0^3} = 0,
    \ese    
    so both amplitudes vanish, which is what we needed. That is we only get non-vanishing 3-point amplitude for complex momenta. 
    \item We can actually obtain the form of $\Cref{eqn:A--+}$ from helicity arguments. We have 
    \bse 
        \begin{split}
            h_1 A(1^-,2^-,3^+) &= h_2 A(1^-,2^-,3^+) = - A(1^-,2^-,3^+), \qand \\
            h_3A(1^-,2^-,3^+) & = + A(1^-,2^-,3^+),
        \end{split}
    \ese 
    where 
    \bse 
        h_i = \frac{1}{2}\bigg(-\l^{\a}_i \frac{\p}{\p \l_i^{\a}} + \widetilde{\l}^{\dot{\a}}_i \frac{\p}{\p \widetilde{\l}^{\dot{\a}}_i} \bigg).
    \ese 
    If we then use the ansantz 
    \bse 
        A(1^-,2^-,3^+) = \la 12 \ra^X \la 23\ra^Y \la 31 \ra^Z
    \ese
    then act with the helicity operators and compare with the above we get the following simultaneous equations 
    \bse 
        -\frac{1}{2}(X+Z) = -1 \qquad -\frac{1}{2}(X+Y) = -1 \qand -\frac{1}{2}(Y+Z) = 1
    \ese 
    which solves to give us 
    \bse 
        X = 3 \qand Y = Z = -1,
    \ese 
    which gives us \Cref{eqn:A--+Final}. Of course we only really know this is correct up to a normalisation. We could then ask "why didn't we consider the ansatz $A(1^-2^-3^+) = [12]^X[23]^Y[31]^Z$?" The answer is if we do a similar calculation for this we can show that the result violates locality so must be excluded.\footnote{This is a problem on the worksheets, so I will not type the answer here.}
    
    This result is particularly nice because it allows to extend the above result to a MHV 3-point amplitude of particles of spin-$s$ as
    \bse 
        A(1^-,2^-,3^+) = \bigg( \frac{\la 12 \ra^4}{\la 12 \ra \la 23 \ra \la 31 \ra} \bigg)^s.
    \ese 
    \item Using the cyclicity of the colour-ordered amplitudes, we have 
    \bse 
        \begin{split}
            A(1^-,2^+,3^-) & = A(3^-,1^-,2^+) = \frac{\la 13 \ra^4}{\la 12 \ra \la 23 \ra \la 31 \ra} \\
            A(1^+,2^-,3^-) & = A(2^-,3^-,1^+) = \frac{\la 23 \ra^4}{\la 12 \ra \la 23 \ra \la 31 \ra} 
        \end{split}
    \ese 
\end{itemize}

\bbox 
    Convince yourself that the 3-point MHV amplitude is totally antisymmetric under exchange of particle labels. That is show that 
    \bse 
        A(2^-,1^-,3^+) = - A(1^-,2^-,3^+) \qand A(3^+,2^-,1^-) = - A(1^-,2^-,3^+). 
    \ese 
    \textit{Hint: This can be shown either using the cyclicity properties above or via the explicit expression in terms of polarisation and momenta.}
\ebox 

\section{$n$-Point Amplitudes: Parke-Taylor Formula}

Now it might seem a bit strange that we are finding the 3-point amplitudes, given that we have shown that all \textit{physical} (i.e. real momenta) 3-point amplitudes must vanish. Well, remarkably, the formula for the 3-point MHV amplitudes has simple generalisation to any $n$-point amplitude, known as the \textit{Parke-Taylor formula}:
\mybox{
    \be 
    \label{eqn:ParkeTaylorMHV}
        A(1^+, 2^+, ... , i^-, ... , j^-, ... , n^+) = \frac{\la ij \ra^4 }{\la 12 \ra \la 23 \ra ... \la (n-1) n \ra \la n 1 \ra}
    \ee 
    and 
    \be 
    \label{eqn:ParkeTaylorMHVBar}
        A(1^-, 2^-, ... , i^+, ... , j^+, ... , n^-) = - \frac{[ij]^4 }{ [12][23]... [(n-1) n][n 1]}
    \ee 
}
\noindent with the $\widetilde{\l}_1 \propto \widetilde{\l}_2$ etc conditions applied as with \Cref{eqn:A--+Final,eqn:A++-Final}.

\br 
\label{rem:MHVClarity}
    Note that in the Parke-Taylor formulas we have two of the polarisations being different, rather then just one (as in \Cref{eqn:AmppmNPoint}). Of course for the 3-point amplitude these two cases coincide (as $3-2=1$). This remark is included just to clarify what is meant by MHV/$\overline{\text{MHV}}$: two polarisations differ from the rest. 
\er 

The Parke-Taylor formulas are the MHV and $\overline{\text{MHV}}$ expressions, but for the 3-point amplitudes we also showed that $A(\pm\pm\pm)=0$, so the natural question is "does this hold for higher point amplitudes?" 

\bcl 
\label{claim:Apm}
    Yes, it does hold. That is 
    \be 
    \label{eqn:ApmNPoint}
        A(\pm \pm ... \pm) = 0.
    \ee 
\ecl 

In order to prove the above claim, we need to introduce the following Lemma.

\bl 
\label{Lem:nPoint3Point}
    A tree-level $n$-point diagram can have, at most, $(n-2)$ 3-point vertices. 
\el 

\bbox 
    Prove \Cref{Lem:nPoint3Point}. \textit{Hint: Prove it inductively.}
\ebox 

\bq
    (Of \Cref{claim:Apm}). Note that 
    \ben[label=(\roman*)]
        \item A gluon amplitude can be expressed as a sum of Feynman diagrams with a polarisation vector associated to each external leg. 
        \item Since the amplitude is a Lorentz scalar, each polarisation vector \textit{must} be contracted with either another polarisation vector or an external momentum. 
        \item If it is contracted to an external momentum, it must be through a 3-point vertex as the 4-point vertex doesn't involve momenta. 
        \item By \Cref{Lem:nPoint3Point}, there are at most $(n-2)$ 3-point vertices, each of which could potentially contract an external polarisation vector with an external momentum. Therefore in each diagram we must have at least two polarisation vectors which aren't contracted with an external momentum, and so must be contracted with each other.
        \item Recall that if we pick the reference momenta $q_i=q_j$ then $\epsilon_i^{\pm}\cdot \epsilon_j^{\pm}=0$. Hence if they all have the same helicity, as in \Cref{eqn:ApmNPoint}, we can pick the reference momenta such that $\epsilon_i\cdot \epsilon_j = 0$ for all $\{i,j\}$, and since each diagram contains at least one such contraction the amplitude must vanish. 
    \een 
\eq 

We now return to what we were trying to say in \Cref{rem:MHVClarity}; using similar arguments to the above proof we can show that 
\mybox{
    \be 
    \label{eqn:AmppmNPoint}
        A(\mp \pm \pm ... \pm ) = 0 \qquad \forall n > 3.
    \ee 
}
\noindent In particular, setting $q_i=p_1$ for all $i>1$ ensures that all inner products of polarisation vectors vanish: from \Cref{eqn:PolarisationInnerProducts} we see
\bse 
    \epsilon_i^{\pm}\cdot \epsilon_j^{\pm} = 0 \qquad \forall i,j > 1 \qquad \text{since } q_i=q_j
\ese 
and 
\bse 
    \epsilon_i^{\pm}\cdot \epsilon_1^{\mp} = 0 \qquad \forall i>1 \qquad \text{since } q_i=p_1.
\ese 

\br 
    Recall that for real momenta \Cref{eqn:AmppmNPoint} also holds for $n=3$, but for complex momenta setting $q_i=p_1$ for all $i>1$ generally results in $0/0$ for 3-point kinematics. In this case we can define non-trivial 3-point MHV/$\overline{\text{MHV}}$ amplitudes.
\er 

Hopefully now \Cref{rem:MHVClarity} is more clear, and we summarise it in the following definition. 

\bd[Maximal Helicity Violating]
    An amplitude is said to be \textit{maximal helicity violating} (MHV)\footnote{Of course we have used a convention to distinguish MHV from $\overline{\text{MHV}}$, but hopefully that is clear by now.} if exactly two of the helicities differ from the rest. 
\ed 

MHV gets its name from the fact that it corresponds to the biggest change of helicity from incoming to outgoing. This might not seem obvious at first but we have to remember that we take all of our particles to be outgoing and that switching them to incoming flips the helicity. So $A(-,-,+,+,...,+)$ corresponds to $n$ \textit{outgoing} particles with 2 negative helicities and $(n-2)$ positive ones. If we change $m<(n-2)$ of the positive helicity particles to be our incoming ones, we get a situation where we have $m$ negative helicity incoming particles going to $2$ negative helicity and $(n-2-m)$ positive ones. If we considered $A(-,-,-,+,...,+)$ we would end up with $m$ negative to $3$ negative and $(n-3-m)$ positive, so the total change in helicity is less. The case with 3 differing helicities is known as \textit{NMHV}, where the "N" stands for "next". Similarly $A(-,-,-,-,+,...,+)$ is called NNMHV, or more simply $\text{N}^2$MHV. In general
\bse 
    A(1^-,2^-,...,k^-,(k+1)^+,...,n^+) \qquad \text{is N}^{k-2}\text{MHV}.
\ese 


Hopefully this explanation is clear, and in order to help, we depict the idea diagrammatically below: the left-hand diagram is MHV while the right-hand side is NMHV.

\begin{center}
    \btik 
        \begin{scope}[xshift=-4cm]
            \draw[thick] (0,0) circle [radius=0.75cm];
            \midarrow[rotate around={-35:(0,0)}] (-0.75,0) -- (-1.75,0);
            \node at (-1.7,1.2) {$-$};
            \midarrow[rotate around={-55:(0,0)}] (-0.75,0) -- (-1.75,0);
            \node at (-1.1,1.6) {$-$};
            \midarrow[rotate around={-75:(0,0)}] (-0.75,0) -- (-1.75,0);
            \node at (-0.5,2) {$+$};
            \midarrow[rotate around={35:(0,0)}] (0.75,0) -- (1.75,0);
            \node at (1.7,1.2) {$+$};
            \midarrow[rotate around={55:(0,0)}] (0.75,0) -- (1.75,0);
            \node at (1.1,1.6) {$+$};
            \node[rotate around={-12.5:(0,0)}] at (0.25,1.25) {$\dots$};
            %
            \midarrow[rotate around={-35:(0,0)}] (1.75,0) -- (0.75,0);
            \node at (1.7,-1.2) {$-$};
            \midarrow[rotate around={-55:(0,0)}] (1.75,0) -- (0.75,0);
            \node at (1.1,-1.6) {$-$};
            \midarrow[rotate around={35:(0,0)}] (-1.75,0) -- (-0.75,0);
            \node at (-1.7,-1.2) {$-$};
            \midarrow[rotate around={55:(0,0)}] (-1.75,0) -- (-0.75,0);
            \node at (-1.1,-1.6) {$-$};
            \midarrow[rotate around={75:(0,0)}] (-1.75,0) -- (-0.75,0);
            \node at (-0.5,-2) {$-$};
            \node[rotate around={12.5:(0,0)}] at (0.25,-1.25) {$\dots$};
            %
            \node at (0,-2.5) {Incoming};
            \node at (0,2.5) {Outgoing};
        \end{scope}
        \begin{scope}[xshift=4cm]
            \draw[thick] (0,0) circle [radius=0.75cm];
            \midarrow[rotate around={-35:(0,0)}] (-0.75,0) -- (-1.75,0);
            \node at (-1.7,1.2) {$-$};
            \midarrow[rotate around={-55:(0,0)}] (-0.75,0) -- (-1.75,0);
            \node at (-1.1,1.6) {$-$};
            \midarrow[rotate around={-75:(0,0)}] (-0.75,0) -- (-1.75,0);
            \node at (-0.5,2) {$-$};
            \midarrow[rotate around={35:(0,0)}] (0.75,0) -- (1.75,0);
            \node at (1.7,1.2) {$+$};
            \midarrow[rotate around={55:(0,0)}] (0.75,0) -- (1.75,0);
            \node at (1.1,1.6) {$+$};
            \node[rotate around={-12.5:(0,0)}] at (0.25,1.25) {$\dots$};
            %
            \midarrow[rotate around={-35:(0,0)}] (1.75,0) -- (0.75,0);
            \node at (1.7,-1.2) {$-$};
            \midarrow[rotate around={-55:(0,0)}] (1.75,0) -- (0.75,0);
            \node at (1.1,-1.6) {$-$};
            \midarrow[rotate around={35:(0,0)}] (-1.75,0) -- (-0.75,0);
            \node at (-1.7,-1.2) {$-$};
            \midarrow[rotate around={55:(0,0)}] (-1.75,0) -- (-0.75,0);
            \node at (-1.1,-1.6) {$-$};
            \midarrow[rotate around={75:(0,0)}] (-1.75,0) -- (-0.75,0);
            \node at (-0.5,-2) {$-$};
            \node[rotate around={12.5:(0,0)}] at (0.25,-1.25) {$\dots$};
            %
            \node at (0,-2.5) {Incoming};
            \node at (0,2.5) {Outgoing};
        \end{scope}
    \etik  
\end{center}

\br 
    As we mentioned in the footnote above, of course we have picked a convention for what we mean by MHV vs $\overline{\text{MHV}}$, and hopefully it is easy to see what the latter corresponds to diagrammatically. 
\er 

\section{4-Point MHV}

Let's now compute the 4-point MHV colour-ordered amplitude $A_4(1^-,2^-,3^+,4^+)$. Recall that this is computed by summing over colour-ordered Feynman diagrams with a fixed cyclic ordering and no cross legs:\footnote{Note that we take all momenta to flow outwards, as per our convention that we have outgoing particles.}
\begin{center}
    \btik 
        \begin{scope}[xshift=-2.5cm]
            \midarrow (-0.5,0) -- (0.5,0) node [midway, below] {$p$};
            \midarrow[rotate around={45:(-0.5,0)}] (-0.5,0) -- (-1.5,0) node [below] {$2^-$};
            \midarrow[rotate around={-45:(-0.5,0)}] (-0.5,0) -- (-1.5,0) node [above] {$1^-$};
            \midarrow[rotate around={-45:(0.5,0)}] (0.5,0) -- (1.5,0) node [below] {$3^+$};
            \midarrow[rotate around={45:(0.5,0)}] (0.5,0) -- (1.5,0) node [above] {$4^+$};
            \node at (2.75,0) {$+$};
        \end{scope}
        \begin{scope}[xshift=2.5cm]
            \midarrow (0,-0.5) -- (0,0.5) node [midway, right] {$p$};
            \midarrow[rotate around={45:(0,0.5)}] (0,0.5) -- (0,1.5) node [above] {$1^-$};
            \midarrow[rotate around={-45:(0,0.5)}] (0,0.5) -- (0,1.5) node [above] {$4^+$};
            \midarrow[rotate around={45:(0,-0.5)}] (0,-0.5) -- (0,-1.5) node [below] {$3^+$};
            \midarrow[rotate around={-45:(0,-0.5)}] (0,-0.5) -- (0,-1.5) node [below] {$2^-$};
            \node at (2.5,0) {$+$};
        \end{scope}
        \begin{scope}[xshift=7.5cm]
            \midarrow[thick, rotate around={45:(0,0)}] (0,0) -- (1,0);
            \midarrow[thick, rotate around={-45:(0,0)}] (0,0) -- (1,0);
            \midarrow[thick, rotate around={45:(0,0)}] (0,0) -- (-1,0);
            \midarrow[thick, rotate around={-45:(0,0)}] (0,0) -- (-1,0);
            \node at (1,1) {$4^+$};
            \node at (-1,1) {$1^-$};
            \node at (-1,-1) {$2^-$};
            \node at (1,-1) {$3^+$};
        \end{scope}
    \etik 
\end{center}

Now recall that the full amplitude is then obtained by dressing $A_4(1^-,2^-,3^+,4^+)$ with the trace $\Tr[T^{a_1}T^{a_2}T^{a_3}T^{a_4}]$ and summing over non-cyclic permutations of external legs. Hence, although the colour-ordered amplitude only has $3$ diagrams, the full amplitude has $3\times 3!=18$ diagrams. This is already a massive simplification, but we can actually further simplify this by a clever choice of reference vectors for the external gluons:
\be
\label{eqn:4PointReferenceMomenta}
    q_1 = q_2 = p_3 \qand q_3 = q_4 = p_2,
\ee    
then all polarisation products vanish expect for 
\be 
\label{eqn:Epsilon14}
    \epsilon_1^- \cdot \epsilon_4^+ = \frac{ \la q_4 1\ra [4q_1] }{\la q_4 4\ra [q_11] } = \frac{\la 21 \ra [43]}{\la 24\ra [31]}.
\ee 
Then the third diagram (the 4-point diagram) vanishes as it is given by 
\bse 
    i (\epsilon_1^-\cdot \epsilon_3^+) (\epsilon_2^-\cdot \epsilon_4^+) = 0.
\ese 

Now consider the t-channel type diagram and look at the bottom 3-point vertex 
\begin{center}
    \btik 
        \node[left] at (-1.25,0) {\Large{$iV_{23p}^{\mu} \,\, =$}};
        \midarrow (0,0) -- (0,1) node [above] {$p,\mu$};
        \midarrow[rotate around={45:(0,0)}] (0,0) -- (-1,0) node [below] {$2^-$};
        \midarrow[rotate around={-45:(0,0)}] (0,0) -- (1,0) node [below] {$3^+$};
    \etik 
\end{center}
This is given by\footnote{We leave the $\pm$ symbols on the polarisations implicit, they can easily be read off from the diagram.}
\bse 
    \begin{split}
        iV_{23p}^{\mu} & = -i\sqrt{2}\big[ (\epsilon_2\cdot \epsilon_3) p_2^{\mu} + \epsilon_3^{\mu} (p_3\cdot \epsilon_2) + \epsilon_2^{\mu} (p\cdot \epsilon_3) \big] \\
        & = -i\sqrt{2}\big[ \epsilon_3^{\mu} (q_2\cdot \epsilon_2) + \epsilon_2^{\mu}(-p_2\cdot \epsilon_3 - p_3\cdot \epsilon_3) \big] \\
        & = -i\sqrt{2} \big[ \epsilon_3 (-q_3\cdot \epsilon_3)\big] \\
        & = 0,
    \end{split}
\ese 
where we have used our reference momenta choice, \Cref{eqn:4PointReferenceMomenta}, $p=-p_2-p_3$ and
\bse 
    \epsilon_2\cdot \epsilon_3 = q_2\cdot \epsilon_2 = q_3 \cdot \epsilon_3 = p_3 \cdot \epsilon_3 = 0.
\ese 
So the second diagram also vanishes and we only need to consider the first one! Using the colour-ordered Feynman rules, we are then just left with 
\bse 
    A_4(1^-,2^-,3^+,4^+) = iV_{12p}^{\mu} i V_{34-p}^{\nu} \frac{-i\eta_{\mu\nu}}{p^2}.
\ese
We then just need to compute the vertex contributions:
\bse 
    \begin{split}
        iV_{12p}^{\mu} & = -i\sqrt{2}\big[ (\epsilon_1\cdot \epsilon_2) p_1^{\mu} + \epsilon_2^{\mu} (p_2\cdot \epsilon_1) + \epsilon_1^{\mu} (p\cdot \epsilon_2)  \big] = -i\sqrt{2}\big[ \epsilon_2^{\mu} (p_2\cdot \epsilon_1) + \epsilon_1^{\mu} (p\cdot \epsilon_2)  \big] \\
        iV_{34-p}^{\nu} & = -i\sqrt{2}\big[ (\epsilon_3\cdot \epsilon_4) p_3^{\nu} + \epsilon_4^{\nu} (p_4\cdot \epsilon_3) - \epsilon_3^{\nu} (p\cdot \epsilon_4)  \big] = -i\sqrt{2}\big[ \epsilon_4^{\nu} (p_4\cdot \epsilon_3) - \epsilon_3^{\nu} (p\cdot \epsilon_4)  \big],
    \end{split}
\ese
where again we have used $\epsilon_1\cdot \epsilon_2 = 0 = \epsilon_3\cdot \epsilon_4$. Now when we do the $\eta_{\mu\nu}$ contraction, and recall that with our reference momenta, the only non-vanishing contraction is $\epsilon_1\cdot\epsilon_4$ we see that only one term survives and we have 
\bse 
    A_4(1^-,2^-,3^+,4^+) = 2 \frac{(p\cdot \epsilon_2)(p_4\cdot \epsilon_3)(\epsilon_1\cdot \epsilon_4)}{p^2}.
\ese    

We now want to express this in terms of our spinor inner products. Firstly, we have 
\bse
    p^2 = (p_1+p_2)^2 = -\la 12 \ra [12].
\ese 
Next we have
\bse 
    p\cdot \epsilon_2 = -p_1\cdot \epsilon_2 - p_2\cdot \epsilon_2 = - p_1\cdot \epsilon_2,
\ese 
and so using\footnote{Putting the helicity labels back on for clarity of where the right-hand sides come from.}
\bse 
    \begin{split}
        p_1\cdot \epsilon_2^- & = \frac{[q_21]\la 12\ra }{\sqrt{2}[2q_2]} = \frac{[31]\la 12\ra }{\sqrt{2}[23]} \\
        p_4\cdot \epsilon_3^+ & = \frac{\la q_34\ra [43]}{\sqrt{2}\la q_33\ra } = \frac{\la 24\ra [43]}{\sqrt{2}\la 23\ra }.
    \end{split}
\ese 
Putting these together with \Cref{eqn:Epsilon14}, we have 
\bse 
    A_4(1^-,2^-,3^+,4^+) = \frac{1}{\la 12\ra [12]} \frac{[31]\la 12\ra }{[23]} \frac{\la 24\ra [43]}{\la 23\ra } \frac{\la 21 \ra [43] }{\la 24\ra [31]}  = \frac{\la 21\ra [43]^2}{[12][23]\la 23 \ra }.
\ese

This is not so nice, and seems to be in contrast with our Parke-Taylor result, \Cref{eqn:ParkeTaylorMHV}.\footnote{Of course what we are trying to do here is demonstrate that the Parke-Taylor formula holds at least for the 4-point MHV.} How do we make it nicer? Well we can cheat a bit by using that we want it to obey \Cref{eqn:ParkeTaylorMHV}, so we know we need to remove the all the square inner products and we want $\la 12 \ra^4$ in the numerator. 

So how do we achieve this? Well first we note that momentum conservation gives us
\bse 
    \la 34 \ra [43] = (p_3+p_4)^2 = (p_1+p_2)^2 = \la 21 \ra [12],
\ese 
so if we multiply and divide by $\la 34\ra$, we get (using $\la 21 \ra^2 = \la 12\ra^2$)
\bse 
    A_4(1^-,2^-,3^+,4^+) = \frac{\la 12 \ra^2 [43]}{[23]\la 23\ra \la 34\ra}. 
\ese 

We now need to remove the other $[43]$ and the $[23]$ in the denominator. The trick to note in order to do this is the content of the next exercise.
\bbox 
    Using momentum conservation, $p_1+p_2=p_3+p_4$, show that 
    \bse 
        [43]\la 41 \ra = - [23] \la 21\ra. 
    \ese 
\ebox  

\noindent Using the result from this exercise, we can multiply and divide by $\la 41 \ra$ to obtain 
\be 
\label{eqn:A4MHVResult}
    A_4(1^-,2^-,3^+,4^+) = \frac{\la 12\ra^3}{\la 23 \ra \la 34\ra \la 41 \ra} = \frac{\la 12\ra^4}{\la 12 \ra \la 23 \ra \la 34\ra \la 41 \ra},
\ee 
which agrees exactly with the Parke-Taylor formula. 

\section{Photon-Decoupling Identity \& Schowten Identity}

Great, we have shown that the Parke-Taylor formula holds for the 4-point amplitude by explicitly calculating $A(1^-,2^-,3^+,4^+)$. We saw it was quite a lot of work in order to obtain this, so it would be extremely useful if we could relate this result to other MHV 4-point amplitudes. That is, we want to use the result above to write down the result for $A_4(1^-,2^+,3^-,4^+)$. 

Again, if we accept the Parke-Taylor identity as true, this is trivial --- simply change the values for $i,j$ in \Cref{eqn:ParkeTaylorMHV} --- however we want to show that this holds explicitly. In order to do that, we introduce the \textit{photon-decoupling identity}
\mybox{
    \be 
    \label{eqn:PhotonDecoupling}
        A(1,2,3,...,n) + A(2,1,3,...,n) + A(2,3,1,...,n) + ... + A(2,3,...,n,1) = 0,
    \ee 
}
\noindent which can be remembered using the mnemonic "migrate the $1$ through". 

\bbox 
    Prove \Cref{eqn:PhotonDecoupling}. \textit{Hint: Use the fact that pure gluon tree amplitudes can be expressed in the following colour-decomposed form}
    \bse 
        \cA_n = g^{n-2} \sum_{\substack{\text{non-cyclic} \\ \text{perms}}} \Tr[T^{a_1} ... T^{a_n}] A(1,...,n),
    \ese 
    \textit{and then let $T^{a_1}=\b1$.}
\ebox 

\br 
    The above exercise is one set on the course so I don't want to expand on the proof any further, but this remark is included to maybe clear up some potential confusion. Setting $T^{a_1}$ to be the identity corresponds to making one of the gluons into a photon, which is why \Cref{eqn:PhotonDecoupling} vanishes: photons and gluons do not couple. This is where the name "photon-decoupling" comes from. Now it might seem strange to then say that the result must also hold when all the particles are gluons (i.e. there is no photon), but what we have to notice is that \Cref{eqn:PhotonDecoupling} is expressed in terms of the colour-ordered amplitudes, which know \textit{nothing} about the colour structure. That is the $A(1,...,n)$ have no way of knowing if the entries are photons or gluons. 
\er 

Ok so using the photon-decoupling identity we have 
\bse 
    \begin{split}
        A(1^-,2^+,3^-,4^+) & = - A(1^-,2^+,4^+,3^-) - A(1^-,4^+,2^+,3^-) \\
        & = - A(3^-,1^-,2^+,4^+) - A(3^-,1^-,4^+,2^+,3^-) \\
        & = -\bigg( \frac{ \la 13\ra^4 }{\la 12 \ra \la 24 \ra \la 43\ra \la 31\ra } + \frac{ \la 13\ra^4 }{\la 14 \ra \la 42 \ra \la 23\ra \la 31\ra } \bigg) \\
        & = \frac{\la 13 \ra^3}{\la 24\ra} \bigg( \frac{1}{\la 12\ra \la 43\ra } - \frac{1}{\la 14\ra \la 23 \ra} \bigg) \\
        & = \frac{\la 13 \ra^3}{\la 24\ra} \bigg( \frac{\la 14\ra \la 23\ra - \la 12\ra \la 43\ra }{\la 12\ra \la 43\ra \la 14\ra \la 23\ra }  \bigg) \\
        & = \frac{\la 13 \ra^3}{\la 24\ra} \bigg( \frac{\la 14\ra \la 23\ra + \la 12\ra \la 34\ra }{\la 12\ra \la 43\ra \la 14\ra \la 23\ra }  \bigg)
    \end{split}
\ese
where the second line follows from the cyclicity of the colour-ordered amplitudes and then we have used our result from above.

This isn't quite what we want and it's not obvious at this point how to make it more Parke-Taylor-like. Indeed we now need another, very useful, identity known as the \textit{Schowten identity}
\mybox{
    \be 
    \label{eqn:Schowten}
        \la ij \ra \l_k + \la jk \ra \l_i + \la ki \ra \l_j = 0.
    \ee 
}

\bbox
    Prove the Schowten identity, \Cref{eqn:Schowten}. \textit{Hint: Note that the $\l$s are 2-component objects, so can think of $\l_i$ and $\l_j$ as a basis for $\l_k$.}\footnote{Again this is a exercise on the course, and I can't add much more of a hint without writing the answer. If any readers are still confused, feel free to drop me an email.}
\ebox 

Using the Schowten identity, we have 
\bse 
    \la 14 \ra \la 23 \ra + \la 12 \ra \la 34 \ra = - \la 13 \ra \la 42 \ra = \la 13 \ra \la 24 \ra,
\ese 
which if we plug in to the expression above gives us 
\bse 
    A(1^-,2^+,3^-,4^+) = \frac{\la 13\ra^4 }{\la 12 \ra \la 23 \ra \la 34\ra \la 41 \ra },
\ese 
which is in agreement with the Parke-Taylor identity. 

So we have managed to prove the Parke-Taylor identity for $n=3$ and $n=4$. Of course this does not prove the Parke-Taylor identity in general. For higher $n$-point functions using Feynman diagrams is not feasible. We now move on to prove a recursion relation that will allow us to prove the Parke-Taylor formula inductively.  
\chapter{BCFW Recursion}

As we just said, we now want to obtain some kind of recursive relation that will allow us to compute higher order amplitudes. As the name of this chapter suggests, such a recursive relation is known as the Britto–Cachazo–Feng–Witten (BCGW) recursion relation. 

So we want to compute a tree-level $n$-point amplitude $A_n$. The method to do this seems slightly strange at first, but hopefully will be clear as we move forward. The idea is to deform legs $1$ and $n$ in such a way that preserves momentum conservation, namely 
\be 
\label{eqn:p1pnDeformation}
    \begin{split}
        p_1 \to \widehat{p}_1(z) &:= p_1 - zq \\
        p_n \to \widehat{p}_n(z) &:= p_n + zq 
    \end{split}
\ee 
where $z\in \C$ and $q$ is some other momentum (i.e. it is a four component object that we can contract with the other momenta $p_i$). It's clear that we can preserved momentum conservation as $\widehat{p}_1+\widehat{p}_n = p_1+p_n$. Now if we want to require $\widehat{p}_1^2 = 0 = \widehat{p}_n^2$ (i.e. they are on-shell momenta) then we must impose 
\be 
\label{eqn:Hatp2=0Condition}
    q^2 = q \cdot p_1 = q \cdot p_n = 0.
\ee 
The next important thing to note is that, because $z\in\C$, the momenta are complex and so we get non-trivial $3$-point amplitudes. This is important as, as we will see, all higher point amplitudes are made up from the $3$-point amplitude (i.e. we have a recursion relation). 

Doing the deformation \Cref{eqn:p1pnDeformation}, we obtain the deformed amplitude $\widehat{A}_n(z)$. The question we now want to ask is "what are its analytical properties?" It is the sum of deformed Feynman diagrams, and so it is a rational function\footnote{A function $f(x)$ is said to be rational if it can be written in the form $f(x) = \frac{P(x)}{Q(x)}$, where $P(x)$ and $Q(x)$ are polynomials in $x$, with $Q(x)$ not being the zero function.} of $z$. Moreover, $\widehat{A}_n(z=0)=A_n$ only has poles when denominators of Feynman propagators become zero, i.e. the virtual exchange particles go on-shell. It follows from this that $\widehat{A}(z)$ only has poles at values of $z$ for which the deformed propagators go on-shell. Near such poles, the amplitude factorises into a product of two on-shell deformed amplitudes, which we denote $\widehat{A}_L(z_i)$ and $\widehat{A}_R(z_i)$ (for "left" and "right"). We depict the idea diagrammatically, in order to try help further explain why this is the case. 

\begin{center}
    \btik 
        \begin{scope}[xshift=-3cm]
            \node at (-3.25,0) {\Large{$\lim$}};
            \node at (-3.25,-0.5) {$z\to z_i$};
            \draw[thick] (0,0) circle [radius=0.75cm];
            \midarrow[rotate around={-45:(0,0)}] (0.75,0) -- (1.75,0) node [right] {$\widehat{n}$};
            \midarrow[rotate around={45:(0,0)}] (0.75,0) -- (1.75,0) node [right] {$5$};
            \midarrow[rotate around={90:(0,0)}] (0.75,0) -- (1.75,0) node [above] {$4$};
            \midarrow[rotate around={-45:(0,0)}] (-0.75,0) -- (-1.75,0) node [left] {$3$};
            \midarrow (-0.75,0) -- (-1.75,0) node [left] {$2$};
            \midarrow[rotate around={45:(0,0)}] (-0.75,0) -- (-1.75,0) node [left] {$\widehat{1}$};
            \node at (1.25,0.1) {\Large{$\vdots$}};
            \node at (0,0) {$\widehat{A}_n(z)$};
        \end{scope}
        \node at (-0.25,0) {\Large{$=$}};
        \begin{scope}[xshift=3cm]
            \draw[thick] (0,0) circle [radius=0.75cm];
            \midarrow (0,-0.75) -- (0,-1.75) node [below] {$\widehat{1}$};
            \midarrow[rotate around={45:(0,0)}] (-0.75,0) -- (-1.75,0) node [left] {$2$};
            \midarrow (-0.75,0) -- (-1.75,0) node [left] {$3$};
            \midarrow (0,0.75) -- (0,1.75) node [above] {$i-1$};
            \node[rotate around={45:(0,0)}] at (-0.9,0.9) {\Large{$\dots$}};
            \node at (0,0) {$\widehat{A}_L(z_i)$};
            %
            \begin{scope}[xshift=3cm]
                \draw[thick] (0,0) circle [radius=0.75cm];
                \midarrow (0,-0.75) -- (0,-1.75) node [below] {$\widehat{n}$};
                \midarrow[rotate around={-45:(0,0)}] (0.75,0) -- (1.75,0) node [right] {$n-1$};
                \midarrow (0.75,0) -- (1.75,0) node [right] {$n-2$};
                \midarrow (0,0.75) -- (0,1.75) node [above] {$i$};
                \node[rotate around={-45:(0,0)}] at (0.9,0.9) {\Large{$\dots$}};
                \node at (0,0) {$\widehat{A}_R(z_i)$};
            \end{scope}
            \midarrow (2.25,0) -- (0.75,0) node [midway, below] {$\widehat{P}_i(z)$};
        \end{scope}
    \etik 
\end{center}
\noindent where 
\be 
\label{eqn:HatPi}
    \widehat{P}_i(z) := (\widehat{p}_1 + p_2 + ... + p_{i-1}) = P_i - zq
\ee 
where we have also defined 
\bse 
    P_i := \sum_{j=1}^{i-1} p_j.
\ese 
It is important to note the notation used: the subscript on the capital $\widehat{P}_i/P_i$ tells us where the sum ends, it is \textit{not} the momentum of the $i$-th leg, which is lowercase $p_i$. $\widehat{P}_i$ is the momentum of the deformed propagator. This is hopefully clear from the sum expression for $P_i$. We will clarify what $z_i$ is in a moment.

\br 
    We should clarify that $i\in \{3,...,n-1\}$,\footnote{As otherwise either $\widehat{A}_L(z_i)$ or $\widehat{A}_R(z_i)$ only contains one particle and so is not a proper amplitude} and so the labelling on the right-hand diagram is simply illustrative (i.e. the leg $3$ could actually belong to $\widehat{A}_R(z_i)$, for example). 
\er 

Ok so why does having our deformed propagator cause the amplitude to split? Well the pedagogical answer goes as follows: when the deformed propagator goes on shell it essentially behaves like an external particle, and so we can view it as such. However it is not actually external (it is not asymptotically free), but instead it connects two other amplitudes with genuine on-shell external legs. In other words we can almost imagine process $\widehat{A}_R(z_i)$ happening, with $\widehat{P}_i(z)$ being an external particle. This particle then becomes an incoming particle in the, otherwise completely separate, amplitude $\widehat{A}_L(z_i)$. 

It is hopefully clear that we must view $\widehat{P}_i(z)$ as an outgoing particle for one amplitude and ingoing for the other, otherwise we would break momentum conservation. Similarly it follows that $\widehat{1}$ and $\widehat{n}$ must be in different subamplitudes (as in the diagram above), as if they appeared in the same subamplitude, the propagator will not be deformed and therefore will not pick up a pole (i.e. will not become "like an external particle"). 

With the idea hopefully cleared up, let's now look at this in a bit more detail. Squaring \Cref{eqn:HatPi}, we get
\bse 
    \widehat{P}_i^2(z) = P_i^2 - 2zP_i \cdot q = -2P_i\cdot q\bigg( z - \frac{P_i^2}{2P_i\cdot q} \bigg).
\ese 
Now we want $\widehat{P}_i^2(z=z_i)=0$ --- i.e. it goes on-shell at $z=z_i$ --- and so we can conclude that
\mybox{ 
    \be 
    \label{eqn:zi}
        z_i := \frac{P_i^2}{2P_i\cdot q}
    \ee 
}

For colour-ordered amplitudes $P_i$ is always given by the sum of adjacent momenta as written above, but more generally $\widehat{P}_i$ just corresponds to the sum of external momenta in $\widehat{A}_L$. This is just the statement of momentum conservation: we have $\widehat{P}_i$ in and the sum $\widehat{p}_1 + ... + p_{i-1}$ out. 

Ok let's look at the subamplitudes in a bit more detail. For each factorisation we must actually sum over all on-shell states propagating between $\widehat{A}_R$ and $\widehat{A}_L$. For gluons this corresponds to a sum over helicities, $s$, which follows from 
\bse 
    \eta_{\mu\nu} = -\sum_{s=\pm1} \epsilon_s^{\mu} \epsilon_{-s}^{\nu}.
\ese 
So we are left with
\be 
\label{eqn:HatAnProduct}
    \begin{split}
        \lim_{z\to z_i} \widehat{A}_n(z) = \frac{1}{z-z_i}\bigg(-\frac{1}{2P_i\cdot q}\bigg) \sum_{s=\pm 1} & \widehat{A}_L\big( \widehat{1}(z_i), 2, ... , i-1, -\widehat{P}_i(z_i)^{-s}\big) \\
        &\times \widehat{A}_R\big( +\widehat{P}_i(z_i)^s, i, ... , n-1, -\widehat{n}(z_i)\big),
    \end{split}
\ee 
where we have used 
\bse 
    \frac{1}{\widehat{P}_i^2} = \frac{1}{z-z_i}\bigg(-\frac{1}{2P_i\cdot q}\bigg).
\ese 
Also note that $\widehat{P}_i$ appears in $\widehat{A}_L$ with a minus sign as it is \textit{ingoing} there and our convention is that all particles are outgoing. For the same reason we put it to the power $-s$, as we flip the helicity when we go from outgoing to incoming. Similarly $\widehat{P}_i$ appears in $\widehat{A}_R$ with a positive sign and a positive power of $s$.

In summary, after deforming the amplitude, the residues of the poles correspond to products of lower point amplitudes, the $\widehat{A}_L$ and $\widehat{A}_R$. By summing over all residues, we can then reconstruct the deformed amplitude from lower-point amplitudes. The original amplitude is then obtain by setting $z=0$. This can be proven using Cauchy's theorem. In particular suppose that 
\be
\label{eqn:limZInftyHatA}
    \lim_{z\to\infty} \widehat{A}_n(z) = 0,
\ee 
then the following contour integral vanishes
\bse 
    \oint_{z=\infty} \frac{dz}{2\pi i} \frac{\widehat{A}_n(z)}{z} = 0,
\ese
where the limit on the integration means we take the contour over an infinite radius circle centered around the origin of the complex $z$ plane. Alternatively, we could define $w := 1/z$ and then take the contour around a small circle around the origin $w=0$, and the integral simply implies that $\widehat{A}_n(w=0)=0$. 

Why is this useful? Well \Cref{eqn:HatAnProduct} tells us that the contour integral also corresponds to the sum over residues of poles inside the contour, notably $z=0$, and the poles where $\widehat{A}_n(z)$ factorises into lower-point amplitudes:
\bse 
    0 = \widehat{A}_n(z=0) + \sum_{i=3}^{n-1} \frac{1}{z_i} \bigg(-\frac{1}{2P_i\cdot q}\bigg) \sum_{s=\pm 1} \widehat{A}_L^{-s}(z_i) \widehat{A}_r^s(z_i),
\ese 
where the sum over $i$ accounts for the different factorisation channels (i.e. the different distribution of the legs between $\widehat{A}_L$ and $\widehat{A}_R$). Then using \Cref{eqn:zi} and $\widehat{A}_n(z=0) = A_n$, we finally conclude 
\mybox{
    \be 
    \label{eqn:BCFWRecursion}
        A_n = \sum_{i=3}^{n-1} \sum_{s=\pm1} \widehat{A}_L^{-s}(z_i) \frac{1}{P_i^2} \widehat{A}_R^s(z_i).
    \ee 
}
\noindent This is the BCFW recursion relation and it tells us that the undeformed amplitude can be computed by summing over products of deformed lower point on-shell amplitudes times undeformed propagators. The deformation parameter of each term corresponds to the value of $z$ for which the deformed propagator goes on-shell. In this way we can recursively compute higher point amplitudes from lower point amplitudes. Note that although the $3$-point amplitude is unphysical itself (it vanishes when we impose our real momenta condition), it is the building block for \textit{all} higher point amplitudes. In other words, the BCFW relation allows us to compute the entire $S$-matrix given only the $3$-point amplitude. 

\section{Comments On Generality Of BCFW}

Although we have focused on colour ordered Yang-Mills amplitudes (i.e. gluons), the BCFW recursion relation can be applied much more generally. Indeed note that the only assumptions we have made is the existence of a Feynman diagram expansion and \Cref{eqn:limZInftyHatA}. This means that, in particular, the BCFW recursion relation can be applied to gravity! This is particularly powerful as the Einstein-Hilbert action has an infinite number of Feynman vertices,\footnote{See the effective field theory section of QFT II to see why.} which makes standard Feynman diagram calculations very complicated. Using the BCFW recursion relation only the $3$-point amplitude needs to be known, and this can be deduced from little-group scaling and locality (as we did above). This last point is very important as it tells us we don't even need to know the Lagrangian of the theory!

Furthermore, the BCFW recursion relation holds in any spacetime dimension $d\geq 4$. In these lectures we have of course focused on $d=4$, as here we can use our spinor techniques. In particular we can define the deformation as follows: 
\bse 
    \l_1 \to \widehat{\l}_1 := \l_1 - z\l_n \qand \widetilde{\l}_n \to \widehat{\widetilde{\l}}_n := \widetilde{\l}_n + z\widetilde{\l}_1,
\ese 
so that 
\bse 
    \begin{split}
        \widehat{p}_1^{\dot{\a}\a}(z) & =  \widehat{\widetilde{\l}}_1^{\dot{\a}} \widehat{\l}_1^{\a} = \widetilde{\l}_1^{\dot{\a}} (\l_1 - z\l_n)^{\a} = p_1^{\dot{\a}\a} - z \widetilde{\l}_1^{\dot{\a}} \l_n^{\a}, \qand \\
        \widehat{p}_n^{\dot{\a}\a}(z) & = \widehat{\widetilde{\l}}_n^{\dot{\a}} \widehat{\l}_n^{\a} = (\widetilde{\l}_n +z \widetilde{\l}_1)^{\dot{\a}} \l_n^{\a} = p_n^{\dot{\a}\a} + z \widetilde{\l}_1^{\dot{\a}}\l_n^{\a},
    \end{split}
\ese
so if we define 
\bse 
    q^{\dot{\a}\a} := \widetilde{\l}_1^{\dot{\a}}\l_n^{\a},
\ese
we have 
\bse 
    \widehat{p}_1 = p_1 - zq \qand \widehat{p}_n = p_n + zq,
\ese 
which is exactly our deformation, \Cref{eqn:p1pnDeformation}. We call this deformation a "$\la 1n]$ shift".

Given this deformation, under what circumstances does \Cref{eqn:limZInftyHatA} hold? To see this, consider deforming the $4$-point MHV amplitude $A_4(1^-,2^-,3^+,4^+)$. Suppose we do a $\la 1^-2^-]$ shift, then from \Cref{eqn:A4MHVResult} we have 
\bse 
    \widehat{A}_4^{--} = \frac{\la \widehat{1} \widehat{2}\ra^3 }{\la \widehat{2}3 \ra \la 34 \ra \la 4\widehat{1}\ra },
\ese 
where the superscript is meant to indicate we are defomring the negative helicity particles. Now note that 
\bse 
    \la \widehat{1}\widehat{2} \ra = \big(\bra{1} - z\bra{2}\big)\ket{2} = \la 12 \ra
\ese 
where we have used $\la 22 \ra =0$. Similarly we have 
\bse 
    \la 4 \widehat{1} \ra = \la 41 \ra - z \la 42 \ra. 
\ese
Then using that $\widehat{2}$ only effects the tilded $\widetilde{\l}_2$ and that the angular $\la ij \ra$ is the inner product w.r.t. untilded $\l_i$ and $\l_j$ we have $\la \widehat{2} 3 \ra = \la 23 \ra$. So in total we have 
\bse 
    \widehat{A}_4^{--} \sim \frac{1}{z},
\ese
which vanishes in the limit $z\to\infty$, and so \Cref{eqn:limZInftyHatA} is obeyed.

Similarly a $\la 3^+ 4^+ ]$ shift will result in 
\bse 
    \widehat{A}_4^{++} \sim \frac{1}{z},
\ese 
and a $\la 4^+ 1^-]$ shift gives 
\bse 
    \widehat{A}^{+-}_4 \sim \frac{1}{z}. 
\ese 
However if we consider a $\la 1^- 4^+]$ shift then we have 
\bse 
    \widehat{A}^{-+}_4 = \frac{\la \widehat{1}2\ra^3}{\la 23 \ra \la 3\widehat{4}\ra \la\widehat{4}\widehat{1}\ra }
\ese 
and 
\bse 
    \la \widehat{1}2 \ra = \la 12 \ra - z\la 42\ra, \qquad \la 3 \widehat{4}\ra = \la 34 \ra \qand \la \widehat{1}\widehat{4} \ra = \la 14 \ra,
\ese 
so in total 
\bse 
    A_4^{-+} \sim z^3,
\ese
which obviously doesn't obey \Cref{eqn:limZInftyHatA}. 

Hence we see that the BCFW recursion relation applies for $\la -- ]$, $\la ++ ]$ and $\la +-]$ shifts but \textit{not} for a $\la -+]$ shift. Although we have only showed this for $n=4$, it turns out this result holds for any $n$ and any MHV degree.\footnote{For a proof see \href{https://arxiv.org/pdf/0801.2385.pdf}{arXiv:0801.2385}.}

\section{Inductive Proof Of Parke-Taylor Formula}

As we said at the end of the last chapter, our motivation for introducing the BCFW recursion relation is so that we can prove the Parke-Taylor identity for $n>4$ point functions by induction. This is exactly what we do now. 

Consider an $n$-point MHV amplitude, $A(1^-,2^+,...,(n-1)^+,n^-)$, and assume that the Parke-Taylor formula holds for all lower point MHV amplitudes. We now perform a $\la 1^- n^-]$ shift. The important thing to notice is that there are only two types of diagrams can contribute, namely:
\begin{center}
    \btik 
        \begin{scope}
            \draw[thick] (0,0) circle [radius=0.75cm];
            \midarrow[rotate around={-45:(0,0)}] (0,-0.75) -- (0,-1.75) node [below] {$\widehat{1}^-$};
            \midarrow[rotate around={45:(0,0)}] (0,0.75) -- (0,1.75) node [above] {$2^+$};
            \node at (0,0) {$\widehat{A}_L$};
            %
            \draw[thick] (3,0) circle [radius=0.75cm];
            \midarrow[rotate around={45:(3,0)}] (3,-0.75) -- (3,-1.75) node [right] {$\widehat{n}^-$};
            \midarrow[rotate around={75:(3,0)}] (3,-0.75) -- (3,-1.75) node [right] {$(n-1)^+$};
            \midarrow[rotate around={-35:(3,0)}] (3,0.75) -- (3,1.75) node [above] {$3^+$};
            \node[rotate around={-70:(0,0)}]  at (4.1,0.4) {\Large{$\dots$}};
            \node at (3,0) {$\widehat{A}_R$};
            \midarrow (2.25,0) -- (0.75,0) node [midway, below] {$\widehat{P}$};
            \node at (1,0.25) {$+$};
            \node at (2,0.25) {$-$};
        \end{scope}
        \begin{scope}[yshift=-5cm]
            \draw[thick] (0,0) circle [radius=0.75cm];
            \midarrow[rotate around={-45:(0,0)}] (0,-0.75) -- (0,-1.75) node [below] {$\widehat{1}^-$};
            \midarrow[rotate around={-75:(0,0)}] (0,-0.75) -- (0,-1.75) node [left] {$2^+$};
            \midarrow[rotate around={35:(0,0)}] (0,0.75) -- (0,1.75) node [above] {$(n-2)^+$};
            \node[rotate around={70:(0,0)}]  at (-1.1,0.4) {\Large{$\dots$}};
            \node at (0,0) {$\widehat{A}_L$};
            %
            \draw[thick] (3,0) circle [radius=0.75cm];
            \midarrow[rotate around={45:(3,0)}] (3,-0.75) -- (3,-1.75) node [right] {$\widehat{n}^-$};
            \midarrow[rotate around={-45:(3,0)}] (3,0.75) -- (3,1.75) node [above] {$(n-1)^+$};
            \node at (3,0) {$\widehat{A}_R$};
            \midarrow (2.25,0) -- (0.75,0) node [midway, below] {$\widehat{P}$};
            \node at (1,0.25) {$-$};
            \node at (2,0.25) {$+$};
        \end{scope}
    \etik 
\end{center}

\noindent The reason we only have these two diagrams is because, as we showed before (\Cref{eqn:AmppmNPoint}), 
\bse 
    A_n(-, +, +, ..., +) = 0 \qquad \forall n >3,
\ese
and every other diagram would contain a subamplitude of exactly that form. 

Now each of the diagrams above contain a $\overline{\text{MHV}_3}$ amplitude (the $\widehat{A}_L$ in the first one and the $\widehat{A}_R$ in the second one). Then recall that for such a amplitude we have all the $\l$s proportional, and since we are shifting $\l_1$ but not $\l_n$ (we shift $\widetilde{\l}_n$), it will be possible to choose $z$ such that the $3$-point kinematics is satisfied by $\widehat{A}_L$ in the first diagram. However, in general the $3$-point kinematics will \textit{not} be satisfied by $\widehat{A}_R$ in the second diagram. This is just because the $\widehat{A}_L$ contains $\l_1$ but $\widehat{A}_R$ contains $\l_n$, and we can shift the former but not the latter. This tells us that the second diagram must vanish for generic kinematics. We can verify this more explicitly. 

For the second diagram we have\footnote{Note we have a minus sign here as $\widehat{P}$ is given by the external momenta of $\widehat{A}_L$, which is $\widehat{p}_1 + p_2 + ... + p_{n-2}$. Then momentum conservation, $\widehat{p}_1 + p_2 + ... + p_{n-1} + \widehat{p}_n=0$ gives us the minus sign.}
\bse
    \widehat{P}(z) = -\big(\widehat{p}_n(z) + p_{n-1}\big) = -\big(p_n + p_{n-1}\big) - zq
\ese 
with $q= \widetilde{\l}_1 \l_n$ from before. Therefore, denoting the pole value as $z_*$, we have 
\bse 
    z_* = \frac{(p_{n-1}+p_n)^2}{-2(p_{n-1}+p_n)\cdot q} = - \frac{ \la (n-1)n \ra [n(n-1)] }{\la (n-1)n\ra [1(n-1)]} = - \frac{[n(n-1)]}{[1(n-1)]}. 
\ese 
From here we have 
\bse 
    \begin{split}
        -\widehat{P}(z_*) & = \widetilde{\l}_{n-1}\l_{n-1} + \widetilde{\l}_n \l_n - \frac{[n(n-1)]}{[1(n-1)]}\widetilde{\l}_1 \l_n \\
        & = \widetilde{\l}_{n-1}\l_{n-1} + \bigg(\widetilde{\l}_n + \frac{[(n-1)n]}{[1(n-1)]}\widetilde{\l}_1\bigg) \l_n \\
        & = \widetilde{\l}_{n-1}\l_{n-1} + \bigg(\frac{\widetilde{\l}_n[1(n-1)] + [(n-1)n]\widetilde{\l}_1}{[1(n-1)]}\bigg) \l_n \\
        & = \widetilde{\l}_{n-1}\l_{n-1} + \frac{[1n]}{[1(n-1)]} \widetilde{\l}_{n-1}\l_n \\
        & = \widetilde{\l}_{n-1}\bigg( \l_{n-1} +\frac{[1n]}{[1(n-1)]}\l_n\bigg),
    \end{split}
\ese 
where we have used the Schowten identity 
\bse 
    [1(n-1)]\widetilde{\l}_n + [(n-1)n]\widetilde{\l}_1 = [1n]\widetilde{\l}_{n-1}.
\ese 
We can write this all in terms of $\l$s using $\widehat{P}= \widetilde{\l}_{\widehat{P}} \l_{\widehat{P}}$:
\be 
\label{eqn:LambdaHatP}
    -\widetilde{\l}_{\widehat{P}} \l_{\widehat{P}} =  \widetilde{\l}_{n-1}\bigg( \l_{n-1} +\frac{[1n]}{[1(n-1)]}\l_n\bigg).
\ee 
Now from \Cref{eqn:A++-Final} we have 
\bse 
    \widehat{A}_R(z_*) = - \frac{[\widehat{P}(n-1)]^3}{[(n-1)\widehat{n}][\widehat{n}\widehat{P}]},
\ese    
but then from \Cref{eqn:LambdaHatP} we have 
\bse 
    [\widehat{P}(n-1)] \sim [(n-1)(n-1)] = 0,
\ese
and 
\bse 
    \begin{split}
        [\widehat{n}\widehat{P}] \sim [\widehat{n}(n-1)] = \big([n| +z_*[1|\big)|n-1] = [n(n-1)] - \frac{[(n-1)n]}{[1(n-1)]}[1(n-1)] = 0
    \end{split}
\ese 
which also gives us $[(n-1)\widehat{n}]=0$, and so 
\bse 
    \widehat{A}_R(z_*) = \frac{0^3}{0^2} = 0,
\ese 
and so the second diagram vanishes, as claimed. 

So we only need to consider the first diagram above\footnote{Hopefully the initial motivating simplifications that arise from using this approach to amplitudes has become clear at this point.} which gives us
\bse 
    A_n(1^-,2^+,...,(n-1)^+,n^-) = - \frac{[2(-\widehat{P})]^3}{[-\widehat{P}\widehat{1}][\widehat{1}2]} \frac{1}{(p_1+p_2)^2} \frac{\la \widehat{n}\widehat{P}\ra^3}{\la \widehat{P}3\ra ... \la (n-1)\widehat{n}\ra}.
\ese
Now let 
\bse 
    \l_{-\widehat{P}} = \l_{\widehat{P}} \qand \widetilde{\l}_{-\widehat{P}} = - \widetilde{\l}_{\widehat{P}} \qquad \implies \qquad \widetilde{\l}_{-\widehat{P}} \l_{-\widehat{P}} = - \widetilde{\l}_{\widehat{P}} \l_{\widehat{P}}.
\ese 
Then using 
\bse 
    |\widehat{P}] \la \widehat{P}| = \widetilde{\l}_{\widehat{P}}^{\dot{\a}} \l_{\widehat{P}}^{\a} = \widehat{P} = \widehat{p}_1+p_2 = | \widehat{p}_1]\la\widehat{p}_1| + |p_2] \la p_2|,
\ese 
we have 
\bse 
    [2(-\widehat{P})]\la \widehat{n}\widehat{P}\ra = [2\widehat{P}]\la \widehat{P}\widehat{n}\ra = [2| \big(\widehat{p}_1]\la\widehat{p}_1| + |p_2] \la p_2|\big) \ket{\widehat{n}} = [21]\la 1 n\ra, 
\ese
where we have also used $\ket{\widehat{n}}=\ket{n}$ and $\la \widehat{1}n\ra = \la 1n\ra$. Similarly we have 
\bse 
    [-\widehat{P}\widehat{1}]\la \widehat{P}3 \ra = [\widehat{1}\widehat{P}] \la \widehat{P} 3\ra = [\widehat{1}| \big( \widehat{p}_1]\la\widehat{p}_1| + |p_2] \la p_2| \big) \ket{3} = [12]\la 23\ra. 
\ese 
Then simply from $[\widehat{1}| = [1|$ and $\ket{\widehat{n}}=\ket{n}$, we also have
\bse 
    [\widehat{1}2] = [12] \qand \la (n-1)\widehat{n}\ra = \la (n-1)n\ra.
\ese 
Finally, we have 
\bse 
    (p_1 + p_2)^2 = \la 12 \ra [21],
\ese 
so in total we have 
\bse 
    A_n(1^-,2^+,...,(n-1)^+,n^-) = \frac{\la n1\ra^3 }{\la 12 \ra \la 23 \ra ... \la (n-1)n\ra },
\ese 
which is the Parke-Taylor formula, \Cref{eqn:ParkeTaylorMHV}.\footnote{Note that in the numerator we have $\la n1\ra$, this accounts for the minus sign missing from cancelling $\la 1n\ra^4/\la n1\ra = -\la 1n\ra^3$ from \Cref{eqn:ParkeTaylorMHV}.}
\chapter{Symmetries Of Amplitudes}

As we have seen, scattering amplitudes have tremendous mathematical simplicity. As we explained at the beginning of the course, very often this is because of underlying symmetries which are sometimes hidden from the point of view of the action and standard Feynman diagram techniques. For the rest of the course, we will explore the conformal symmetries of tree-level Yang-Mills amplitudes and their SUSY extension. We will then introduce twistors and briefly describe how they can be used to realise Yangian symmetry of $\cN=4$ SYM and formulate a worldsheet description. 

\section{Review Of Conformal Group}

As we have an entire course on CFT, this is just a brief section to recap the relevant information needed here. 

The generators of the conformal group are 
\ben[label=(\roman*)] 
    \item Poincar\'{e}
        \ben 
            \item Translations:
            \bse 
                P_{\mu} = -i \frac{\p}{\p x^{\mu}}.
            \ese 
            \item Lorentz:
            \bse 
                M_{\mu\nu} = i\bigg(x_{\mu}\frac{\p}{\p x^{\nu}} - x_{\nu}\frac{\p}{\p x^{\mu}}\bigg).
            \ese 
        \een 
    \item Dilatations:
    \bse 
        D = -ix^{\mu}\frac{\p}{\p x^{\mu}}.
    \ese
    \item Special Conformal:
    \bse 
        K_{\mu} = i\bigg( x^2 \frac{\p}{\p x^{\mu}} - 2x_{\mu} x^{\nu}\frac{\p}{\p x^{\nu}} \bigg)
    \ese 
    The special conformal transformations are related to the translations by the inversion operator,
    \bse 
        I(x^{\mu}) = \frac{x^{\mu}}{x^2},
    \ese 
    simply as $K_{\mu} = IP_{\mu}I$.
\een 

\section{$4D$ SYM}

In 4D Minkowski spacetime, these generate $SO(2,4)$. $4D$ SYM theory has no dimensionful parameters, and therefore enjoys classical conformal symmetry, which is broken quantum mechanically. Quantum mechanics enters at loop level, so we expect the tree-level amplitudes to be conformally invariant for $4D$ SYM. 

\subsection{Generators In Spinor Form}

How is this symmetry realised for amplitudes? Well, in principal, this can be deduced by Fourier transforming the generators to momentum space and changing to spinor variables by the prescriptions above. However very little is gained from this calculation, and so here we simply write our the generators in spinor form and verify that they annihilate MHV amplitudes. Given the spinor form, it is not difficult to verify that they obey the conformal algebra relations (i.e. the commutator relations). 

\ben[label=(\roman*)] 
    \item Translations: 
    \be 
    \label{eqn:Translations}
        p^{\dot{\a}\a} = \widetilde{\l}^{\dot{\a}}\l^{\a}.
    \ee 
    \item Lorentz:
    \be 
    \label{eqn:Lorentz}
        m_{\a\beta} = \l_{(\a} \frac{\p}{\p \l^{\beta)}} \qand \widetilde{m}_{\dot{\a}\dot{\beta}} = \widetilde{\l}_{(\dot{\a}} \frac{\p}{\p \widetilde{\l}^{\dot{\beta})}},
    \ee 
    where the brackets indicate index symmetrisation. 
    \item Dilatations:
    \be 
    \label{eqn:Dilatations}
        d = \frac{1}{2} \l^{\a} \frac{\p}{\p \l^{\a}} + \widetilde{\l}^{\dot{\a}}\frac{\p}{\p \widetilde{\l}^{\dot{\a}}} + \b1. 
    \ee 
    \item Special Conformal:
    \be 
    \label{eqn:SpecialConformal}
        k_{\a\dot{\a}} = \frac{\p}{\p \l^{\a} \widetilde{\l}^{\dot{\a}}}.
    \ee 
\een 

Let's verify that they annihilate the tree-level MHV amplitudes 
\bse 
    A_n^{\text{MHV}} \equiv A(1^+,..., i^-,...,j^-,...,n^+) = \del^4\bigg(\sum_{i=1}^n p_i \bigg) \frac{\la ij \ra}{\la 12\ra ... \la n 1\ra }
\ese 
where the delta function imposes momentum conservation. We put this in as we can then treat the spinor variables in the rest of the expression as independent variables. We will use
\bse 
    p := \sum_{i=1}^n p_i
\ese 
to lighten notation. The full symmetry generators are obtained by defining the above generators for each external leg and then summing them all up. 

\ben[label=(\roman*)]
    \item Translations: these are trivial as we simply get
    \bse 
        \bigg(\sum_{i=1}^n p_i\bigg) \del^4(p) = p \del^4(p) = 0.
    \ese    
    \item Lorentz: this follows from the invariance of $\la jk \ra$ and $[jk]$, 
    \bse 
        \begin{split}
            \sum_{i=1}^n m_{i,\a\beta} \la jk \ra & = \sum_{i=1}^n \l_{i(\a} \frac{\p}{\p \l_i^{\beta)}} \l_j^{\g}\l_{k\g} \\
            & = \frac{1}{2}\big(\l_{j\a} \del^{\g}_{\beta} \l_{k\g} + \l_{j}^{\g} \l_{k\a} \epsilon_{\g\beta}\big) + (\a \leftrightarrow \beta) \\
            & = \frac{1}{2}\big( \l_{j\a} \l_{k\beta} - \l_{j\beta} \l_{k\a}\big) + (\a \leftrightarrow \beta) \\
            & = 0,
        \end{split}
    \ese 
    and then similarly for the $\sum\widetilde{m}_{i,\dot{\a}\dot{\beta}}[jk]$ calculation.
    \item Dilations: First note that 
    \bse 
        \sum_{i=1}^n\frac{1}{2} \l^{\a} \frac{\p}{\p \l^{\a}} + \widetilde{\l}^{\dot{\a}}\frac{\p}{\p \widetilde{\l}^{\dot{\a}}} 
    \ese 
    just counts mass dimension. So from 
    \bse 
        \big[\del^4(p)\big] = -4, \qquad \big[\la ij\ra \big] = 4 \qand \big[ \big(\la 12 \ra ... \la n1 \ra\big)^{-1} \big] = -n
    \ese 
    and 
    \bse 
        \sum_{i=1}^n \b1 = n,
    \ese 
    we have 
    \bse 
        d A_n^{\text{MHV}} = (-4+4-n)A_n^{\text{MHV}} + nA_n^{\text{MHV}} = 0. 
    \ese 
    \item Special Conformal: this take a bit more work. First let's introduct that notation 
    \bse 
        \widetilde{A}_n^{\text{MHV}} := \frac{\la ij \ra}{\la 12\ra ... \la n 1\ra } \qquad \implies \qquad A_n^{\text{MHV}} = \del^4(p) \widetilde{A}_n^{\text{MHV}}.
    \ese
    Then, noting 
    \bse 
        \frac{\p \widetilde{A}_n^{\text{MHV}}}{\p \widetilde{\l}_i^{\dot{\a}}} = 0,
    \ese 
    we have 
    \bse 
        \begin{split}
            \sum_{i=1}^n k_{i,\a\dot{\a}} A_n^{\text{MHV}} & = \sum_{i=1}^n \frac{\p}{\p \l_i^{\a} \widetilde{\l}_i^{\dot{\a}}} \Big[ \del^4(p) \widetilde{A}_n^{\text{MHV}}\Big] \\
            & = \sum_{i=1}^n \frac{\p}{\p \l_i^{\a}} \bigg[ \frac{\p p^{\dot{\beta}\beta}}{\p \widetilde{\l}_i^{\dot{\a}}} \frac{\p \del^4(p)}{\p p^{\dot{\beta}\beta}} \widetilde{A}_n^{\text{MHV}}\bigg] \\
            & = \sum_{i=1}^n \frac{\p}{\p \l_i^{\a}} \bigg[ \l_i^{\beta} \del^{\dot{\beta}}_{\dot{\a}} \frac{\p \del^4(p)}{\p p^{\dot{\beta}\beta}} \widetilde{A}_n^{\text{MHV}}\bigg] \\
            & = \sum_{i=1}^n \bigg[ \del^{\beta}_{\a} \del^{\dot{\beta}}_{\dot{\a}} \frac{\p \del^4(p)}{\p p^{\dot{\beta}\beta}} \widetilde{A}_n^{\text{MHV}} + \l_i^{\beta} \del^{\dot{\beta}}_{\dot{\a}} \widetilde{\l}_i^{\dot{\g}} \del^{\g}_{\a} \frac{\p^2 \del^4(p)}{\p p^{\dot{\g}\g} \p p^{\dot{\beta}\beta}} \widetilde{A}_n^{\text{MHV}} + \l_i^{\beta} \del^{\dot{\beta}}_{\dot{\a}} \frac{\p \del^4(p)}{\p p^{\dot{\beta}\beta}} \frac{\p \widetilde{A}_n^{\text{MHV}}}{\p \l_i^{\a}} \bigg] \\
            & = \bigg[ n \del^{\beta}_{\a} \del^{\dot{\beta}}_{\dot{\a}} \frac{\p \del^4(p)}{\p p^{\dot{\beta}\beta}}  + p^{\dot{\beta}\beta} \frac{\p^2 \del^4(p)}{\p p^{\dot{\a}\beta} \p p^{\dot{\beta}\a}}\bigg] \widetilde{A}_n^{\text{MHV}} + \frac{\p \del^4(p)}{\p p^{\dot{\a}\beta}} \sum_{i=1}^n \l_i^{\beta} \frac{\p \widetilde{A}_n^{\text{MHV}}}{\p \l_i^{\a}}.
        \end{split}
    \ese 
    This looks like a horrible mess, but now we note that 
    \bse 
        \l_{\beta} \frac{\p}{\p \l^{\a}} = \l_{(\beta} \frac{\p}{\p \l^{\a)}} + \l_{[\beta} \frac{\p}{\p \l^{\a]}} = m_{\a\beta} + \frac{1}{2} \epsilon_{\beta\a} \l^{\g} \frac{\p}{\p \l^{\g}}, 
    \ese 
    so the last term in the expression above simply gives 
    \bse 
        \frac{\p \del^4(p)}{\p p^{\dot{\a}\beta}} \sum_{i=1}^n \l_i^{\beta} \frac{\p \widetilde{A}_n^{\text{MHV}}}{\p \l_i^{\a}} = \frac{\p \del^4(p)}{\p p^{\dot{a}\beta}} \sum_{i=1}^n \bigg( {m^{\beta}}_{\a} + \frac{1}{2}\del^{\beta}_{\a} \l_i^{\g}\frac{\p}{\p \l_i^{\g}}\bigg) A_n^{\text{MHV}} = (4-n) \frac{\p \del^4(p)}{\p p^{\dot{\a}\beta}} A_n^{\text{MHV}},
    \ese 
    where we have used that the Lorentz generators annihilate $A_n^{\text{MHV}}$ and that the derivative term will just give the weight, $(4-n)$, as in the dilatation calculation. 
    
    So in total we are left with 
    \bse 
        \sum_{i=1}^n k_{i,\dot{\a}\a} A_n^{\text{MHV}} = \bigg[ 4 \frac{\p \del^4(p)}{\p p^{\dot{\a}\a}}  + p^{\dot{\beta}\beta} \frac{\p^2 \del^4(p)}{\p p^{\dot{\beta}\a} \p p^{\dot{\a}\beta}}\bigg]\widetilde{A}_n^{\text{MHV}}.
    \ese
    The claim is now that 
    \bse 
        p^{\dot{\beta}\beta} \frac{\p^2 \del^4(p)}{\p p^{\dot{\beta}\a} \p p^{\dot{\a}\beta}} = -4 \frac{\p \del^4(p)}{\p p^{\dot{\a}\a}}.
    \ese 
    We can verify this by integrating against a test function:
    \bse 
        \begin{split}
            \int d^4p \, F(p) p^{\dot{\beta}\beta} \frac{\p^2 \del^4(p)}{\p p^{\dot{\beta}\a} \p p^{\dot{\a}\beta}} & = - \int d^4p \, \bigg( \frac{\p F}{\p p^{\dot{\beta}\a}} p^{\dot{\beta}\beta} + F 2\del^{\beta}_{\a}\bigg) \frac{\p \del^4(p)}{\p p^{\dot{\a}\beta}} \\
            & = \int d^4p \, \bigg( \frac{\p^2 F}{\p p^{\dot{\a}\beta} \p p^{\dot{\beta}\a}}p^{\dot{\beta}\beta} + \frac{\p F}{\p p^{\dot{\beta}\a}} 2\del^{\dot{\beta}}_{\dot{\a}} + \frac{\p F}{\p p^{\dot{\a}\beta}} 2\del^{\beta}_{\a} \bigg)\del^4(p) \\
            & = 4 \int d^4p \, \frac{\p F}{\p p^{\dot{\a}\a}} \del^4(p) \\
            & = -4 \int d^4p \, F(p) \frac{\p \del^4(p)}{\p p^{\dot{\a}\a}},
        \end{split}
    \ese 
    and so we have 
    \bse 
        \sum_{i=1}^n k_{i,\a\dot{\a}} A_n^{\text{MHV}} = 0. 
    \ese
\een 

So we have shown that our given conformal symmetry generators annihilate the MHV tree-level amplitudes. 

\subsection{SUSY}

It's possible to extend the conformal group by introducing Grassman-odd variables, $\eta^A$, and defining the new generators
\be 
\label{eqn:SUSYqsDefinitions}
    \begin{split}
        q^{\a A} & := \l^{\a} \eta^A, \\
        \widetilde{q}^{\dot{\a}}_{A} &:= \widetilde{\l}^{\dot{\a}} \frac{\p}{\p\eta^A}, \\
        s_{\a A} & := \frac{\p}{\p \l^{\a}} \frac{\p }{\p\eta^A} \\
        \widetilde{s}_{\dot{\a}}^A & := \eta^{A} \frac{\p}{\p \widetilde{\l}^{\dot{\a}}} 
    \end{split}
\ee 
These generators obey 
\be 
\label{eqn:qsAnticommutators}
    \{ q^{\a A}, \widetilde{q}^{\dot{\a}}_B \} = \del^A_B p^{\dot{\a}\a} \qand \{ s_{\a A}, \widetilde{s}_{\dot{\a}}^{B} \} = \del^B_A k_{\a\dot{\a}} 
\ee 

\bbox 
    Verify \Cref{eqn:qsAnticommutators} hold.  
\ebox 

The range of $A$ corresponds to the amount of SUSY.\footnote{See the SUSY course for more on this.} For example if we want to describe $\cN=4$ SYM, $A=1,2,3,4$, which is the maximally symmetric YM theory in 4$D$.

In contrast to the fact that conformal symmetry is broken quantum mechanically in pure YM, superconformal symmetry in $\cN=4$ SYM persists in the quantum theory. Due to its high degree of symmetry, many quantities in $\cN=4$ SYM can be analytically computer to high, and in some cases arbitrary, order in coupling. In that sense, $\cN=4$ SYM can be thought of as a toy model for QCD. For this reason we will focus on the case of $\cN=4$ SYM for the remainder of this course. 

In addition to the SUSY generators, \Cref{eqn:SUSYqsDefinitions}, the superconformal group also has $R$-symmetry generators, 
\be 
\label{eqn:RGenerator}
    {r^A}_B := \eta^{A} \frac{\p}{\p \eta^B} - \frac{1}{4} \del^A_B \eta^C \frac{\p}{\p \eta^C},
\ee 
which generate $SU(4)$ rotations in the $\eta^A$ space. We summarise the the superconformal generators in the following table, also listing their mass dimension (or equivalently their dilatation weight)

\mybox{
    \begin{center}
    	\begin{tabular}{@{} C{4cm} C{4cm} C{4cm} @{}}
    		\toprule
    		Name & Symbol & Mass Dimension \\ 
    		\midrule
    		Translations & $p^{\dot{\a}\a}$ & $1$ \\ \\
    		SuperPoincar\'{e} & $q^{\a A}$ and $\widetilde{q}^{\dot{\a}}_A$ & $1/2$ \\ \\
    		Lorentz & $m_{\a\beta}$ and $\widetilde{m}_{\dot{\a}\dot{\beta}}$ & $0$ \\ \\
    		Dilatations & $d$ & $0$ \\ \\
    		R & ${r^A}_B$ & 0 \\ \\
    		SuperConformal & $s_{\a A}$ and $\widetilde{s}_{\dot{\a}}^A$ & $-1/2$ \\ \\
    		Special Conformal & $k_{\a\dot{\a}}$ & $-1$ \\
    		\bottomrule
    	\end{tabular}
    \end{center}
}

The reason we listed their mass dimensions is that it gives us insight into the structure of the superconformal symmetry algebra, namely through the fact that the weight of the commutator of objects of weight $w_1$ and $w_2$ is $(w_1+w_2)$. In particular this allows us to see that we can obtain all other generators from commuting $q,\widetilde{q}, s$ and $\widetilde{s}$.\footnote{Note this is just a heuristic argument. In itself it does not tell us how to relate $p^{\dot{\a}\a}$, say, to the $q$s and $s$s. The  point is that you can anticipate the structure of the superalgebra simply from dimensional analysis and index structure, but you would need to carry out a calculation to prove it. The bottom line is that if you want to verify that an amplitude is superconformal, it's sufficient to show that its annihilated by the $q$s and $s$s.}

We also modify the definition of the helicity operator to be 
\be 
\label{eqn:HelicityModified}
    h := \frac{1}{2}\bigg( - \l^{\a} \frac{\p}{\p \l^{\a}} + \widetilde{\l}^{\dot{\a}}\frac{\p}{\p \widetilde{\l}^{\dot{\a}}} + \eta^A \frac{\p}{\p \eta^A}\bigg).
\ee 
We can easily check that this commutes with all of the generators, and so it represents a central extension of the algebra. In other words it can be used to define a central charge $c$. For example, we find that 
\bse 
    \{ q^{\a A}, s_{\beta B}\} = {m^{\a}}_{\beta} \del^A_B + \del^{\a}_{\beta} {r^A}_B + \frac{1}{2}\del^{\a}_{\beta} \del^A_B \big(d+c\big),
\ese 
where 
\mybox{
    \be 
    \label{eqn:CentralCharge}
        c = 1-h.
    \ee 
}
\noindent We the define \textit{superamplitudes} to be annihilated by all superconformal generators and the central charge $c$, just as our `normal' amplitudes were annihilated by the conformal generators. It follows from \Cref{eqn:CentralCharge}, then, that \textit{all} superamplitudes have helicity $+1$:
\bse  
    c A = 0 \qquad \iff \qquad hA = A.
\ese 

What exactly \textit{is} a superamplitude? It should be thought of as the scattering amplitude for superfields, which encode all on-shell degrees of freedom.\footnote{See SUSY course for more details on superfields.} For $\cN=4$ SYM with gauge group $SU(N)$ the superfields are given by the component decomposition 
\bse 
    \Phi(\l,\widetilde{\l},\eta) = g_+(p) + \eta^A \psi_A(p) + \frac{1}{2}\eta^A\eta^B \phi_{AB}(p) + \frac{1}{3!} \eta^A\eta^B\eta^C \epsilon_{ABCD}\bar{\psi}^D(p) + \frac{1}{4!}\eta^A\eta^B\eta^C\eta^D \epsilon_{ABCD}g_-(p),
\ese 
where $p = \widetilde{\l}\l$ and where 
\begin{itemize}
    \item $g_{\pm}$ are gluons, i.e. $h=\pm1$, respectively,
    \item $\psi_A$ and $\bar{\psi}^A$ are $4+4$ Fermions, i.e. $h=\pm 1/2$, respectively
    \item $\phi_{AB} = -\phi_{BA}$ are $6$ scalars, i.e. $h=0$,
\end{itemize}
are known as the \textit{components}, and they all transform in the adjoint representation of $SU(N)$ (we have suppressed the colour indices above). From this it is easy to check that $h\Phi =\Phi$. 

Note that we have $2+6=8=4+4$ Bosons and Fermions, respectively, which is in agreement with the fact that SUSY requires there to be the same number of each. Indeed, SUSY mixes Bosons and Fermions and the SUSY transformations of component fields can be read off from\footnote{For more details on how to do this, again see the SUSY course.}
\bse 
    \del_q \Phi \equiv \xi_{\a A} q^{\a A}\Phi, \qand \del_{\bar{q}} \Phi \equiv \bar{\xi}_{\dot{\a}}^A \bar{q}^{\dot{\a}}_A \Phi.
\ese 

\subsection{Superamplitudes}

A superamplitude encodes gluon amplitudes and all other component amplitudes related by SUSY. In general the superamplitude will be a polynomial in $\eta$s, and the component amplitudes are coefficients of the polynomial.\footnote{Again this idea should be clear from the SUSY course.} Since each component field is associated with a certain monomial of $\eta$s in the superfield, each component amplitude is multiplied by the product of $\eta$ monomials for each component field in the amplitude. 

For example, an $n$-point MHV amplitude $A_n(-,-,+,...,+)$ will be multiplied by $(\eta_1)^4(\eta_2)^4$, where 
\bse 
    (\eta_i)^4 = \frac{1}{4!} \epsilon_{ABCD}\eta_i^A\eta_i^B\eta_i^C\eta_i^D,
\ese 
since each negative helicity gluon, $g_-(p)$, appears with $(\eta)^4$ in the superfield, but the positive helicity gluon, $g_+(p)$, doesn't appear with any $\eta$s.

From this, we see that the superamplitude which contains an MHV amplitude must have Fermionic degree of weight $8$. On the other hand, it must be annihilated by the multiplicative charges
\bse 
    q^{\a A} = \sum_{i=1}^n \l^{\a}_i \eta_i^A \qand p^{\dot{\a}\a} = \sum_{i=1}^n \widetilde{\l}^{\dot{\a}}_i \l^{\a}_i.
\ese
Hence, the MHV superamplitude must take the form 
\bse
    \mathbb{A}_n^{\text{MHV}} = \del^4(p) \del^8(q) P_n(\l_i, \widetilde{\l}_i)
\ese 
where 
\be 
\label{eqn:del8q}
    \del^8(q) := \frac{1}{2^4} \prod_{A=1}^4 q^{\a A}q_{\a}^A = \prod_{A=1}^4 \sum_{i<j} \la ij \ra \eta_i^A \eta_j^A.
\ee 
Moreover, we can determine the function $P_n$ by demanding that the $(\eta_1)^4(\eta_2)^4$ component of the superamplitude is $A_n(-,-,+,...,+)$. Noting that 
\bse 
    \int d^4 \eta_1 d^4 \eta_2 \del^8(q) = \la 12 \ra^4,
\ese 
we see that $P_n$ must be 
\bse 
    P_n = \frac{1}{\la 12\ra ... \la n1 \ra}. 
\ese 
Hence 
\mybox{
    \be 
    \label{eqn:MHVSuperamplitude}
        \mathbb{A}_n^{\text{MHV}} = \frac{\del^4(p)\del^8(q)}{\la 12 \ra ... \la n1 \ra}.
    \ee 
}

\br
    Note that this amplitude encodes all $n$-point gluonic MHV amplitudes. Indeed 
    \bse 
        \int d^4\eta_i d^4\eta_j \mathbb{A}_n^{\text{MHV}} = \frac{\la ij \ra^4}{\la 12 \ra ... \la n1\ra} \del^4(p).
    \ese 
    It also encodes amplitudes with scalars and Fermions. 
\er 

\subsection{Supertwistor Space \& Yangian Symmetry}

Now note that, apart from through $\del^4(p)$, \Cref{eqn:MHVSuperamplitude} does not depend on $\widetilde{\l}$s at all! This has important consequence. In particular let's Fourier transform the amplitude in $\widetilde{\l}_i$ and $\eta_i$, using $\mu_i$ and $\chi_i$ being the Fourier conjugates, respectively: 
\bse 
    \begin{split}
        \mathbb{A}_n(\l_i, \mu_i, \chi_i) & = \int \prod_{i=1}^n \frac{d^2\widetilde{\l}_i d^4\eta_i}{(2\pi)^4} \exp\bigg[ i \bigg( \sum_{i=1}^n \big(\mu_i^{\dot{\a}} \widetilde{\l}_{i\dot{\a}} + \chi_i^A\eta_{iA}\big) \bigg) \bigg] \\
        & \qquad \times \int d^4x d^8\Theta \exp\bigg[ i\bigg( x_{\a\dot{\a}} \sum_{i=1}^n \widetilde{\l}_i^{\dot{\a}} \l_i^{\a} + \Theta_{\a}^A \sum_{i=1}^n \eta_{iA}\l_i^A\bigg) \bigg] \frac{1}{\la 12 \ra ... \la n1 \ra} \\
        & = \int d^4x d^8\Theta \prod_{i=1}^n \del^2\big(\mu_i + x\cdot \l_i\big) \del^4\big(\chi_i + \Theta\cdot \l_i\big) \frac{1}{\la 12 \ra ... \la n1 \ra},
    \end{split}
\ese 
where we have used 
\bse 
    \del^4(p) \del^8(q) = \int d^4x d^8\Theta \exp\bigg[ i\bigg( x_{\a\dot{\a}} \sum_{i=1}^n \widetilde{\l}_i^{\dot{\a}} \l_i^{\a} + \Theta_{\a}^A \sum_{i=1}^n \eta_{iA}\l_i^A\bigg) \bigg].
\ese 
The space $(\l,\mu,\chi)$ is known is \textit{supertwistor space}. We then find that MHV superamplitudes are supported on degree 1 curves in twistor space. $(x,\Theta)$ are the moduli of the curves. More generally, $N^k$MHV amplitudes are supported on $(k+1)$ degree curves. This property of amplitudes can be made manifest by reformulating $\cN=4$ SYM as a string theory whose target space is twistor space. For obvious reasons we do not discuss this further here but further details can be found via \href{https://arxiv.org/pdf/hep-th/0312171.pdf}{hep-th/0312171}.

Twistors are also useful for studying the symmetries of scattering amplitudes. In particular, if we replace 
\bse 
    \widetilde{\l} \to \frac{\p}{\p\mu}, \qquad \frac{\p}{\p \widetilde{\l}} \to \mu, \qquad \eta_A \to \frac{\p}{\p \chi^A}, \qand \frac{\p}{\p \eta_A} \to \chi^A 
\ese
in the definition of the super conformal generators defined earlier, we find that they all become first order differential operators 
\bse 
    {J^{(0)a}}_b = \sum_{i=1}^n Z_i^a \frac{\p}{\p Z_i^b},
\ese 
where 
\bse 
    Z_i^a := \begin{pmatrix}
        \l^{\a} \\
        \mu_{\dot{\a}} \\
        \chi^A 
    \end{pmatrix}.
\ese 
For example, 
\bse 
    p_i = \sum_{i=1}^n \widetilde{\l}_i^{\dot{\a}} \l_i^{\a} \to \sum_{i=1}^n \l_i^{\a} \frac{\p}{\p \mu_{i\dot{\a}}}. 
\ese
Hence, superconformal symmetry is linearly realised in twistor space; twistors transform in the fundamanetal representation of the superconformal group. 

Furthermore, the following non-local operators are also symmetries of the planar amplitudes, which are hidden from the point of view of the action:
\bse 
    {J^{(1)a}}_b = \sum_{i<j} \bigg[ Z_i^a \frac{\p}{\p Z_i^c} Z^c_j \frac{\p}{\p Z_j^b} - (i\leftrightarrow j) \bigg].
\ese  
Note that the algebra of these generators does not close and implies an infinite-dimensional symmetry known as \textit{Yangian symmetry}. This is a hallmark of integrability, and is another hint that $\cN=4$ SYM has a worldsheet description since integrability is usually restricted to $2d$ models. In fact the Yangian symmetry of $\cN=4$ SYM can be understood from IIB string theory on $AdS_5\times S^5$, which is dual to $\cN=4$ SYM at strong coupling. 

In summary, scattering amplitudes are not only essential for relating theory to experiment, they also provide a window into the underlying mathematical structure of QFT. 

% -------------------------------------------------------------------
% Bibliography/Further Readings
% -------------------------------------------------------------------

\chapter*{Useful Texts \& Further Readings}

\section*{Name of section}
\begin{itemize}
    \item Scattering Amps In Gauge Theories, Henn and Plefka 
    \item QFT Srednicki 
    \item Scattering Amps, Eluang and Huarg
\end{itemize}




%\bibliographystyle{agsm} 
%\bibliography{mybibliography} 
%\printbibliography[heading=bibintoc]


% -------------------------------------------------------------------
% Appendices
% -------------------------------------------------------------------

%\begin{appendices}
%\chapter{An Appendix of Some Kind}

\lipsum[1-3]  % Replace with your text

%\end{appendices}

\end{document}
